\chapter{Théorie des ensembles ordonnés}
\label{chp.ordres}

\minitoc

\lettrine{D}{ans} ce chapitre, nous aborderons la théorie des ensembles
ordonnés. Celle-ci est fortement liée à la logique, et permet entre autre de
formaliser des structures logiques comme les algèbres de Boole ou les
algèbres de Heyting.

Les ensembles ordonnés peuvent être vus comme des cas dégénérés de catégories.
Si nous ne donnons pas dans cet ouvrage de rudiments de théorie des catégories,
l'étudiant catégoricien pourra voir dans plusieurs définitions des éléments
similaires à ce qu'il connait en théorie des catégories.

Nous commencerons par définir les idées générales sur les ensembles ordonnés.
Cette première partie introduira les notions élémentaires d'ordre, de borne
supérieure, inférieure, ainsi que les ordres bien fondés se reposant sur les
éléments minimaux.
Nous aborderons ensuite la théorie des treillis, en allant des demi-treillis aux
algèbres de Boole et de Heyting.
Cela motivera notre étude des filtres et ultrafiltres, dont le comportement dans
une algèbre de Boole est fortement simplifiée, et nous aurons alors
l'occasion de démontrer le théorème d'extension d'une théorie cohérente en une
théorie complète.

\section{Ensembles ordonnés}

\noindent
Commençons par donner les définitions les plus élémentaires.

\subsection{Définitions}

\begin{definition}[Ensemble ordonné]
  Un ensemble ordonné est un couple $(X,\leq)$ où $\leq$ est une relation
  binaire sur $X$ vérifiant~:
  \begin{itemize}
  \item $\leq$ est reflexive~: pour tout $x\in X$, $x\leq x$.
  \item $\leq$ est antisymétrique~: si deux éléments $x,y\in X$ sont tels que
    $x\leq y$ et $y\leq x$ alors $x = y$.
  \item $\leq$ est transitive~: si trois éléments $x,y,z\in X$ sont tels que
    $x\leq y$ et $y\leq z$ alors $x\leq z$.
  \end{itemize}

  Un couple $(X,\leq)$ tel que $\leq$ est seulement réflexive et transitive est
  appelé un ensemble pré-ordonné, et $\leq$ est appelé un pré-ordre sur $X$.
\end{definition}

\begin{example}[Ensemble des parties]
  Soit $X$ un ensemble, l'ensemble $(\powerset(X),\subseteq)$ est un ensemble
  ordonné.
\end{example}

La plupart des résultats que nous énoncerons pourront aussi bien se traduire sur
des pré-ordres, en remplaçant en général $x = y$ par $x\leq y \land y \leq x$.
En fait, un pré-ordre peut se ramener à un ordre par un quotient.

\begin{property}
  Soit $(X,\leq)$ un ensemble pré-ordonné, alors la relation
  $\mathord{\sim}\subseteq X \times X$ définie par
  \[x\sim y \defeq x\leq y \land y \leq x\]
  est une relation d'équivalence, et $\leq$ induit un ordre sur
  $\quot{X}{\sim}$.
\end{property}

\begin{proof}
  Vérifions que $\sim$ est une relation d'équivalence~:
  \begin{itemize}
  \item $\sim$ est réflexive~:
    $x\leq x$ et $x\leq x$, donc $x\sim x$.
  \item $\sim$ est symétrique~:
    supposons que $x\sim y$, alors $x\leq y$ et
    $y\leq x$, donc $y\leq x$ et $x\leq y$, donc $y\sim x$.
  \item $\sim$ est transitive~:
    supposons que $x\sim y$ et $y\sim z$, alors
    $x\leq y$ et $y\leq z$, donc $x\leq z$, et de même $z\leq y$ et $y\leq x$
    donc $z\leq z$, ainsi $x\sim z$.
  \end{itemize}

  Pour vérifier que $\leq$ induit un ordre sur $\quot{X}{\sim}$,
  il suffit de vérifier
  que si $x\sim x'$ et $y\sim y'$, alors $x\leq y$ si et seulement si
  $x'\leq y'$.
  Si $x\leq y$ alors comme $x'\leq x$ et $y\leq y'$, on en déduit que
  $x'\leq y'$
  et réciproquement, comme $x\leq x'$ et $y'\leq y$, il vient que
  $x'\leq y'\implies x\leq y$.
\end{proof}

\begin{exercise}
  Soit une signature $\Sign$, montrer que la relation
  $\mathord{\vdP}\subseteq\Formula(\Sign)\times\Formula(\Sign)$ définie comme
  restriction de la relation $\vdP$ que nous avons vue précédemment, mais où
  la liste de gauche ne contient qu'une proposition, est un pré-ordre.
\end{exercise}

\begin{remark}
  La relation $\dvvd$ s'écrit aussi $\equiv$, ce qui est cohérent avec
  notre définition sémantique de $\equiv$, étant donné que l'on sait que
  $\vdP$ et $\vD$ coïncident.

  L'ensemble ordonné construit par $\quot{\Prop(\Sigma)}{\dvvd}$ est appelé
  l'algèbre de Lindenbaum-Tarski (cette construction a été
  introduite dans \cite{LindenbaumTarski}). Nous la verrons plus en détail plus
  tard.
\end{remark}

Une autre définition possible d'un ensemble ordonné est ce que l'on appelle
habituellement un ensemble strictement ordonné. Nous en donnons la définition et
montrons que l'on peut faire correspondre à chaque ordre un ordre strict (et
réciproquement).

\begin{definition}[Ordre strict]
  Un ensemble strictement ordonné est une paire $(X,<)$ où $<$ est une relation
  binaire sur $X$ vérifiant~:
  \begin{itemize}
  \item $<$ est antiréflexive~: pour tout $x\in X, x\not< x$.
  \item $<$ est transitive.
  \end{itemize}
\end{definition}

\begin{proposition}
  On a une bijection entre les ensembles ordonnés et les ensembles strictement
  ordonnés en associant à $(X,\leq)$ l'ensemble strictement ordonné $(X,<)$
  défini par
  \[x < y \defeq x \leq y \land x \neq y\]
\end{proposition}

\begin{proof}
  On définit, dans l'autre sens, la relation $\leq$ à partir de $<$ par
  \[x\leq y \defeq x < y \lor x = y\]
  Les deux faits suivants se vérifient directement~:
  \begin{itemize}
  \item  $((x \leq y)\land x \neq y) \lor x = y$ est équivalent à $x\leq y$~;
  \item $(x < y\lor x = y) \land x \neq y$ est équivalent à $x < y$.
  \end{itemize}

  Montrons que $<$ défini à partir de $\leq$ est un ordre strict~:
  \begin{itemize}
  \item par définition, $x < y\implies x \neq y$, donc $x\not< x$.
  \item si $x < y$ et $y < z$, alors $x \leq z$ par transitivité de $\leq$.
    Si $x = z$ alors $z < y$ par substitution de $x$ par $z$, donc par
    transitivité $y < y$, ce qui est absurde. Donc $x \neq z$.
  \end{itemize}

  Montrons que $\leq$ défini à partir de $<$ est un ordre~:
  \begin{itemize}
  \item par définition, si $x = y$ alors $x\leq y$.
  \item si $x\leq y$ et $y\leq x$, alors soit $x = y$, soit $x < y$. Dans le
    premier cas, on a le résultat voulu. Dans le deuxième cas, on utilise le
    fait que $y\leq x$ pour en déduire que soit $y = x$, soit $y < x$. Encore
    une fois, le cas $y = x$ donne directement ce qu'il faut démontrer.
    Si $y < x$, alors par transitivité avec $x < y$ on en déduit que $x < x$,
    ce qui est absurde.
  \item si $x \leq y$ et $y \leq z$, alors soit $x = y$, auquel cas on voit
    directement que $x\leq z$, soit $x < y$. En utilisant le fait que
    $y \leq z$, on en déduit que soit $y = z$, soit $y < z$. Dans le premier
    cas, il est évident que $x \leq z$. Dans le deuxième cas, par transitivité,
    on en déduit que $x < z$.\qedhere
  \end{itemize}
\end{proof}

On définira donc des propriétés indifféremment sur un ordre ou sur un ordre
strict (si on dit qu'une propriété $\varFF(\mathord{<}))$ définie sur les ordres
stricts est vérifiée pour un ordre $\leq$, cela signifie que l'ordre strict $<$
associé vérifie $\varFF(<)$).

Intéressons-nous d'abord aux éléments~: comme un ordre est (sauf indication du
contraire) partiel, deux éléments peuvent ne pas être mis en relation. Dans ces
cas-là, deux notions distinctes apparaissent~: la comparabilité et la
compatibilité. Pour deux éléments, être comparable est une condition forte,
mais parfois justement trop forte~: la compatibilité signifie juste que deux
éléments peuvent être rejoints.

\begin{definition}[Comparabilité, compatibilité]
  Soit $(X,\leq)$ un ensemble ordonné. On dit que deux éléments $x$ et $y$ sont
  comparables si $x\leq y$ ou $y\leq x$. Deux éléments $x$ et $y$ sont dits
  compatibles vers le haut (respectivement vers le bas) s'il existe $z \in X$
  tel que $x\leq z$ et $y\leq z$ (respectivement $z \leq x$ et $z \leq y$). Deux
  éléments sont dits incompatibles vers le haut (respectivement vers le bas)
  s'ils ne sont pas compatibles vers le haut (respectivement vers le bas).

  Si tous les éléments de $X$ sont comparables, on dit que $X$ est un ensemble
  totalement ordonné.
  Si tous les éléments de $X$ sont compatibles vers le haut (respectivement vers
  le bas), on dit que $X$ est un ensemble ordonné filtrant vers le haut
  (respectivement vers le bas).
\end{definition}

\begin{example}
  Deux éléments de $\mathcal P(X)$ sont toujours compatibles vers le bas en
  prenant comme élément $z$ leur intersection. Pourtant, $\{x\}$ et $\{y\}$ pour
  $x\neq y$ ne sont pas comparables.
\end{example}

\subsection{Dualité}

Relevons dès maintenant un phénomène important en théorie des ordres, appelé la
dualité~: lorsque l'on a un ordre $\leq$, on peut définir la relation $\geq$
par $y\geq x \iff x\leq y$, et cette relation est aussi une relation d'ordre,
inversant quel élément est le plus grand et quel élément est le plus petit.
Ainsi lorsque l'on définit une notion sur un ensemble ordonné $(X,\leq)$, une
notion duale est directement définie en considérant l'objet dans $(X,\geq)$.

Ce phénomène permet de simplifier beaucoup de preuves~: si l'on veut prouver
deux propositions $\varFF$ et $\varFP$ portant sur tous les ensembles ordonnés,
où $\varFF$ est obtenue en remplaçant $\leq$ dans $\varFF$ par $\geq$ (et toutes
les définitions par leurs définitions duales), il suffit de prouver $\varFF$
pour en déduire $\varFP$, qui est une conséquence de $\varFF$ pour l'ordre dual.

Nous allons donc souvent écrire deux propositions pour une seule preuve, car la
proposition duale est souvent aussi importante que la proposition de base.

\subsection{Bornes et majorations}

On peut voir la théorie des ordres sous deux prismes~: le prisme algébrique et
le prisme analytique.

Suivant le prisme algébrique, nous étudions des structures
munies d'une relation (on peut voir cela comme de la théorie des modèles), et
l'accent est alors mis sur les propositions du premier ordre. Par exemple,
la comparabilité ou la compatibilité sont des notions algébriques. De même on
peut parler de la borne supérieure de deux éléments, ou d'être le maximum entre
$x$ et $y$, voire d'un ensemble fini d'éléments.

Suivant le prisme analytique, nous nous intéressons aux parties de l'ensemble
ordonné (potentiellement infinie). Nous parlons donc par exemple de la borne
supérieure d'un ensemble quelconque. Les structures pour lesquelles on peut
appliquer le point de vue analytique sont donc moins nombreuses.
Etudions les définitions basiques de ces principes de bornes supérieures,
de majorant \latinexpr{etc}.

\begin{definition}[Majorant, minorant]
  Soit $(X,\leq)$ un ensemble ordonné et $\varPt\subseteq X$ une partie de $X$.
  On dit que $x$ est un majorant de $\varPt$ si
  \[\forall y\in \varPt, y \leq x\]
  et que $x$ est un minorant de $\varPt$ si
  \[\forall y\in \varPt, x \leq y\]

  Par abus de notation, on écrira $\varPt\leq x$ (respectivmeent $x\leq \varPt$)
  pour dire que $x$ est un majorant (respectivement un minorant) de $\varPt$.
\end{definition}

\begin{definition}[Borne supérieure, borne inférieure]
  Soit $(X,\leq)$ un ensemble ordonné et $\varPt\subseteq X$ une partie de $X$.
  On dit que $x$ est la borne supérieure de $\varPt$ si
  \[\varPt\leq x \land (\forall y\in X, \varPt\leq y \implies x \leq y)\]
  et que $x$ est la borne inférieure de $\varPt$ si
  \[x\leq \varPt \land (\forall y\in X, x\leq \varPt \implies y \leq x)\]
\end{definition}

\begin{proof}
  Puisque nous avons dit \og la\fg{} il est nécessaire de montrer l'unicité de
  la borne supérieure (l'unicité de la borne inférieure se prouve par dualité).

  Supposons que $x,y$ sont des bornes supérieures de $\varPt$. Comme
  $\varPt\leq x$ et $y$ est une borne supérieure, $y\leq x$. De même, comme
  $\varPt\leq y$ et $x$ est une borne supérieure, $x\leq y$. Donc $x = y$ par
  antisymétrie.
\end{proof}

\begin{notation}
  Si la borne supérieure de $\varPt$ existe, on la notera $\opJn \varPt$. Si la
  borne inférieure de $\varPt$ existe, on la notera $\opMt \varPt$.
\end{notation}

\begin{exercise}
  Soit $(X,\leq)$ un ensemble ordonné et $\varPt\subseteq \powerset(X)$ une
  partie telle que $\bigcup \varPt$ admet une borne supérieure et tout élément
  de $\varPt$ admet une borne supérieure, montrer alors que
  \[\opJn \left(\bigcup \varPt\right) =
  \opJn\compr{\opJn a}{a \in \varPt}\]
\end{exercise}

\begin{definition}[Maximum, minimum]
  Soit $(X,\leq)$ un ensemble ordonné et $\varPt\subseteq X$. On dit que $x$ est
  le maximum de $\varPt$ si $\varPt\leq x$ et $x\in \varPt$. On dit que $x$ est
  le minimum de $\varPt$ si $x\leq \varPt$ et $x\in \varPt$.
\end{definition}

\begin{notation}
  On notera $\max \varPt$ et $\min \varPt$ le maximum (respectivement le
  minimum) de $\varPt$.
\end{notation}

\begin{property}
  S'il existe, le maximum (respectivement le minimum) d'une partie est sa borne
  supérieure (respectivement inférieure).
\end{property}

\begin{proof}
  Si $x$ est un majorant de $\varPt$, alors comme $\max \varPt \in \varPt$, on
  en déduit que $\max \varPt \leq x$, d'où le résultat.
\end{proof}

\begin{definition}[\'Elément maximal, minimal]
  Soit un ensemble ordonné $(X,\leq)$ et une partie $\varPt\subseteq X$. On dit
  que $x \in \varPt$ est un élément maximal (respectivement minimal) dans
  $\varPt$ s'il n'existe pas $y\in \varPt$ tel que $x < y$ (respectivement tel
  que $y < x$).
\end{definition}

\begin{exercise}
  Montrer que le maximum d'une partie est un élément maximal de cette partie.
\end{exercise}

Une autre partie intéressante est l'étude de parties sur lesquelles l'ordre
induit se comporte particulièrement bien. C'est ce que nous allons voir avec les
chaînes, les antichaînes et les parties filtrantes.

\begin{definition}[Chaîne, antichaîne]
  Soit $(X,\leq)$ un ensemble ordonné. On dit que $C\subseteq X$ est une chaîne
  de $X$ si l'ensemble $(C,\leq)$ (avec l'ordre induit) est un ensemble
  totalement ordonné. On dit que $A\subseteq X$ est une antichaîne de $X$ si
  l'ensemble $(A,\leq)$ (avec l'ordre induit) ne possède pas deux éléments
  comparables, c'est-à-dire si $(A,\leq)$ est l'ordre discret. On dit que
  $A\subseteq X$ est une antichaîne forte vers le haut (respectivement vers le
  bas) de $X$ si pour tous $x,y\in A$, $x$ et $y$ ne sont pas compatibles vers
  le haut (respectivement vers le bas).
\end{definition}

Cela nous mène à la notion d'ensemble inductif, importante pour établir le lemme
de Zorn, qui est une version de l'axiome du choix plus maniable. Nous
prouverons ce résultat lorsque nous étudierons la théorie des ensembles, mais
pouvons déjà en donner l'énoncé. La preuve est donnée dans la \cref{sbsct.Zorn}.

\begin{definition}[Ensemble inductif]
  Un ensemble ordonné $(X,\leq)$ est dit inductif si toute chaîne admet un
  majorant.
\end{definition}

\begin{remark}
  On peut remarquer que le nom \og ensemble inductif\fg est aussi utilisé dans
  le \cref{chp.induction}. Dans la littérature, la notion d'ensemble inductif
  se réfère principalement à un ensemble ordonné tel que décrit plus haut, mais
  ce nom reste le plus evocateur pour décrire les ensembles définis dans le
  \cref{chp.induction}.
\end{remark}

\begin{theorem}[Lemme de Zorn \cite{ZornLemma}]\label{thm.Zorn}
  Si $(X,\leq)$ est un ensemble inductif et $x\in X$, alors il existe un élément
  maximal $y \in X$ qui est supérieur à $x$.
\end{theorem}

\begin{remark}
  Il existe plusieurs versions de la définition d'un ordre inductif. Par
  exemple, on peut considérer que toute chaîne admet une borne supérieure, et
  non simplement un majorant, ou que ce ne sont pas simplement les chaînes mais
  les chaînes bien ordonnées (dont nous parlons dans la sous-section suivante)
  qui admettent un majorant (respectivement une borne supérieure). Dans tous les
  cas, la version que nous donnons ici est la plus simple, et les versions
  alternatives se prouvent de la même manière.
\end{remark}

Si le nom de lemme de Zorn a été donné, par exemple dans \cite{Tukey}, ce
résultat apparaît déjà avant, chez Kazimierz \name{Kuratowski} dès 1922. Le nom
de Max \name{Zorn} est cependant resté (bien qu'en Russie on l'appelle plutôt
lemme de Kuratowski-Zorn).
L'article \cite{ZornLemma}, lui, présente avant tout l'efficacité de ce
théorème pour simplifier des constructions algébriques dont les preuves se
ressemblaient toutes.

\subsection{Ordre bien fondé et bon ordre}

Nous donnons maintenant les notions d'ordre bien fondé et de bon ordre. Les bons
ordres auront de l'importance pour la théorie des ensembles, et les ordres bien
fondés sont utiles pour étendre le théorème d'induction.

\begin{definition}[Ordre bien fondé]\label{def.wf.ord}
  Soit $(X,\leq)$ un ensemble ordonné. On dit que $\leq$ est bien fondé si toute
  partie $\varPt\subseteq X$ non vide possède un élément minimal, c'est-à-dire
  si
  \[\forall \varPt \subseteq X, \varPt\neq \vide \implies \exists x\in \varPt,
  \forall y\in \varPt, (y\leq x \implies x = y)\]
\end{definition}

\begin{remark}
  Une définition équivalente, sur les ordres stricts, est que $<$ est bien fondé
  si~:
  \[\forall \varPt \subseteq X, \varPt\neq \vide \implies \exists x \in \varPt,
  \forall y \in \varPt, y\nless x\]
\end{remark}

\begin{definition}[Bon ordre]
  Un ensemble ordonné $(X,\leq)$ est un bon ordre lorsque toute partie non vide
  $\varPt\subseteq X$ possède un minimum.
\end{definition}

Une première conséquence évidente est qu'un bon ordre est un ordre bien fondé.
De plus, un bon ordre est total, puisque pour $x,y\in X$ il existe un minimum
à $\{x,y\}$. Ceci caractérise en fait les bons ordres parmi les ordres bien
fondés.

\begin{proposition}
  Si $(X,\leq)$ est un ordre bien fondé total, alors c'est un bon ordre.
\end{proposition}

\begin{proof}
  Supposons que $(X,\leq)$ est un ordre bien fondé total. Soit
  $\varPt\subseteq X$ non vide, comme $\leq$ est un ordre bien fondé, on trouve
  $m$ un élément minimal de $\varPt$. Soit $x\in \varPt$, alors soit $x \leq m$
  soit $m\leq x$. Dans le premier cas, par minimalité de $m$, on en déduit que
  $m = x$, donc que $m\leq x$. Ainsi, dans tous les cas, $m$ est un minorant de
  $\varPt$~: $\varPt$ admet donc un minimum.
  Donc $(X,\leq)$ est un bon ordre.
\end{proof}

Nous avons vu que la notion de structure inductive nous offre le principe
d'induction, qui dans le cas particulier de $\bN$ correspond au principe
de récurrence. Le principe de récurrence forte, lui, est de nature légèrement
différente~: il se base sur la notion d'ordre. Il dit que si pour tout
$k\in\bN$, on peut prouver
\[(\forall i < k, \varFF(i)) \implies \varFF(k)\]
alors on peut en déduire que $\forall n, \varFF(n)$.

Ce procédé peut se généraliser à un ensemble bien fondé~: on veut montrer que
$\varFF(x)$ est vrai pour tout $x$ à partir du fait que
\[(\forall y < x, \varFF(y)) \implies \varFF(x)\]

L'idée derrière cette démonstration est que pour un ordre bien fondé, on a un
moyen d'ordonner les éléments de sorte qu'on peut prouver $\varFF$ élément par
élément en s'appuyant sur le fait que $\varFF$ est vérifié aux étapes
précédentes.

\begin{theorem}[Induction bien fondée]
  Soit $(X,\leq)$ un ensemble ordonné bien fondé. Soit $\varPt\subseteq X$ une
  partie vérifiant
  \[\forall x, ((\forall y < x, y\in \varPt) \implies x\in \varPt)\]
  Alors $\varPt$ est égale à $X$.
\end{theorem}

\begin{proof}
  Supposons que $\varPt\neq X$. On trouve alors $x\in X \setminus \varPt$
  minimal pour $<$. Comme $x$ est minimal pour $<$ dans $X\setminus \varPt$, on
  en déduit que pour tout $y < x$, $y \notin X\setminus \varPt$, c'est-à-dire
  $y\in \varPt$. Par hypothèse, puisque $\forall y < x, y\in \varPt$, on en
  déduit que $x\in \varPt$~: c'est absurde. Ainsi, par l'absurde, on a prouvé
  que $X\setminus \varPt = \vide$, donc que $\varPt = X$.
\end{proof}

\begin{remark}
  En construisant un ensemble inductif à partir d'une signature $(C,\arite)$ en
  prenant une suite d'ensembles $(X_i)_{i\in \bN}$, on peut définir un ordre
  canonique bien fondé, l'ordre de sous-terme, par récurrence sur $i$~:
  \begin{itemize}
  \item pour $X_0$, on a facilement un ordre bien fondé sur $\vide$.
  \item supposons qu'on possède un ordre bien fondé sur $X_n$, alors on définit
    notre ordre sur $X_{n+1}$ par $x_i < c(x_1,\ldots,x_n)$ pour tout
    $i \in \intervIntB{n}$, et en en prenant la clôture transitive.
  \end{itemize}

  Cet ordre, en particulier, permet de généraliser l'hypothèse d'induction de
  notre définition d'ensemble inductif en une induction forte.
\end{remark}

\begin{exercise}
  Montrer que l'ordre défini plus tôt est bien un ordre bien fondé.
\end{exercise}

\noindent
Cela nous permet aussi de généraliser la récursion en la récursion bien fondée.

\begin{notation}
  Pour gagner en lisibilité, si $(X,<)$ est un ensemble strictement ordonné et
  $x \in X$, on notera
  \[\downcl{x} \defeq \compr{y \in X}{y < x}\]
\end{notation}

\begin{theorem}[Récursion bien fondée]\label{thm.recur.bf}
  Soient $X,Y$ deux ensembles, et $<$ un ordre bien fondé sur $X$. Soit
  une fonction
  \[\makeFunName{f}{\powerset(Y)}{Y}\]
  alors il existe une unique fonction $\makeFunName{g}{X}{Y}$ telle que
  \[\forall x \in X, g(x) = f(\compr{g(x_0)}{x_0 \in \downcl{x}})\]
\end{theorem}

\begin{proof}
  On prouve par induction bien fondée que si la fonction est définie pour les
  éléments $y < x$, alors elle est définie de façon unique en $x$. Mais cela
  est directement le cas puisque l'image de $x$ est l'image par une fonction
  (donc définie de façon unique) d'un ensemble déjà défini de façon unique par
  hypothèse d'induction bien fondée, donc $g(x)$ est définie de façon unique.
\end{proof}

On peut donc définir une fonction $\makeFunName{g}{X}{Y}$ avec $(X,<)$ bien
fondé en la définissant en $x$ grâce à ses valeurs pour tous les $y < x$. On
verra dans le \cref{chp.ordinaux} une version plus générale de ce principe dans
le cas des ensembles bien ordonnés.
On aurait pu renforcer ce théorème en utilisant comme
condition que $g(x) = f(g\restr{\downcl{x}})$, permettant non seulement
d'utiliser les $g(y)$ mais aussi les $y$ associés.

\section{Treillis}

Nous allons maintenant étudier le cas des treillis, qui sont des ensembles
ordonnés munis de propriétés de clôture. Pour être exhaustif et situer
exactement quelles conditions mènent à quels résultats, nous allons définir nos
structures de la plus faible à la plus forte, en montrant tous les résultats que
nous souhaitons avoir sur une structure avant de passer à la suivante.

\subsection{Demi-treillis}

Commençons par étudier les demi-treillis~: ceux-ci sont des ensembles ordonnés
avec une direction privilégiée, selon laquelle il existe toujours une borne
pour un ensemble fini.

\begin{definition}[Demi-treillis]
  Un ensemble $(X,\leq)$ est un sup demi-treillis si toute partie $F\subfin X$
  admet une borne supérieure $\opJn F$. C'est un inf demi-treillis si toute
  partie $F\subfin X$ admet une borne inférieure $\opMt F$.
\end{definition}

\begin{property}
  De manière équivalente, un ensemble $(X,\leq)$ est un sup demi-treillis si et
  seulement s'il admet un élément $\opB$ et une opération
  $\makeFunName{\opjn}{X \times X}{X}$ tels que $\opB$ est un minorant de $X$ et
  $x\opjn y$ est la borne supérieure de $x$ et $y$.

  C'est un inf demi-treillis si et seulement s'il admet un majorant $\opT$ et
  une opération $\makeFunName{\opmt}{X \times X}{X}$ prenant la borne inférieure
  de deux éléments.
\end{property}

\begin{proof}
  Si $(X,\leq)$ est un sup demi-treillis, on trouve $\opB$ en prenant
  $\opJn\vide$ et $x\opjn y$ en prenant $\opJn\{x,y\}$.

  Réciproquement, montrons par récurrence sur le cardinal de $F\subfin X$ qu'il
  existe une borne supérieure à $F$~:
  \begin{itemize}
  \item si $F = \vide$, alors $\opJn F = \opB$.
  \item si $F = F' \cup \{x\}$ et $\opJn F'$ existe, alors
    $\opJn F = \opJn \{\opJn F',\opJn \{x\}\} = \opJn F' \opjn x$
    donc $\opJn F$ existe. \qedhere
  \end{itemize}
\end{proof}

Cela nous donne une caractérisation plus algébrique d'un demi-treillis~: c'est
un ensemble muni d'opérations. Cependant, ces opérations ont encore des
définitions trop proches de la théorie des ordres. Nous allons voir que l'on
peut transformer cette définition en une définition purement algébrique,
c'est-à-dire n'utilisant que des opérations finies et des égalités.

\begin{definition}[Monoïde commutatif idempotent]\label{def.monoid.ci}
  Un ensemble $\varMMo$ muni d'une opération $\opM$ et d'un élément $\neuMM$ est
  un monoïde idempotent commutatif si les propriétés suivantes sont vérifiées~:
  
  \vspace{0.2cm}
  \begin{itemize}
    \begin{minipage}{0.5\textwidth}
    \item $\forall x\in \varMMo, x\opM \neuMM = \neuMM \opM x = x$
    \item $\forall x,y,z\in\varMMo, x \opM (y \opM z) = (x \opM y) \opM z$
    \end{minipage}
    \begin{minipage}{0.4\textwidth}
    \item $\forall x,y\in\varMMo, x\opM y = y \opM x$
    \item $\forall x\in \varMMo, x\opM x = x$
    \end{minipage}
  \end{itemize}

  On définit sur un monoïde commutatif idempotent $(\varMMo,\opM,\neuMM)$ les
  relations $\leq_\lor$ et $\leq_\land$ par
  \begin{align*}
    x \leq_{\opjn} y &\defeq x \opM y = y \\
    x \leq_{\opmt} y &\defeq x \opM y = x
  \end{align*}
\end{definition}

\begin{proposition}
  L'ensemble $(X,\opB,\opjn)$ est un sup demi-treillis pour $\leq_{\opjn}$ si et
  seulement si $(X,\opB,\mathord{\opjn})$ est un monoïde commutatif idempotent.

  L'ensemble $(X,\opT,\opmt)$ est inf demi-treillis pour $\leq_{\opmt}$ si et
  seulement si $(X,\opT,\mathord{\opmt})$ est un monoïde commutatif idempotent.
\end{proposition}

\begin{proof}
  Montrons d'abord que $\leq_{\opjn}$ est bien un ordre pour un monoïde
  commutatif idempotent~:
  \begin{itemize}
  \item par idempotence, $x\opjn x = x$ donc $x\leq_{\opjn} x$.
  \item si $x\leq_{\opjn} y$ et $y\leq_{\opjn} x$, alors 
    \begin{align*}
      x &= y \opjn x &\text{car } y \leq_{\opjn} x\\
      &= x \opjn y &\text{par commutativité}\\
      &= y & \text{car } x \leq_{\opjn} y
    \end{align*}
  \item si $x\leq_{\opjn} y$ et $y\leq_{\opjn} z$, alors
    \begin{align*}
      x\opjn z &= x \opjn (y \opjn z) &\text{car } y \leq_{\opjn} z\\
      &= (x \opjn y) \opjn z &\text{par associativité}\\
      &= y \opjn z &\text{car } x \leq_{\opjn} y\\
      &= z &\text{car } y \leq_{\opjn} z
    \end{align*}
  \end{itemize}

  Supposons maintenant que $(X,\opB,\opjn)$ est un sup demi-treillis, alors tout
  d'abord $x\leq y$ si et seulement si $y = \opJn\{x,y\}$, si et seulement si
  $y = x \opjn y$, d'où le fait que $\leq_{\opjn}$ coïncide avec $\leq$. Montrons
  maintenant que $(X,\opB,\opjn)$ est un monoïde commutatif idempotent~:
  \begin{itemize}
  \item $\opB$ est un élément neutre pour $\opjn$~: pour tout $x\in X$,
    $\opB\leq x$ donc $\opB \opjn x = x$, et de même pour $x\opjn\opB$.
  \item $\opjn$ est associatif~:
    $x\opjn (y \opjn z) = \opJn\{x,y,z\} = (x\opjn y)\opjn z$.
  \item $\opjn$ est commutatif~: $x\opjn y = \opJn\{x,y\}=y\opjn x$.
  \item $\opjn$ est idempotent~: $x\opjn x = \opJn\{x\}=x$.
  \end{itemize}

  Réciproquement, montrons que si $(X,\opB,\opjn)$ est un monoïde commutatif
  idempotent, alors $\opB$ est un minorant de $X$ pour $\leq_{\opjn}$ et
  $\opjn$ est la borne supérieure de deux éléments pour $\leq_{\opjn}$~:
  \begin{itemize}
  \item pour tout $x\in X$, $\opB \opjn x = x$ donc $\opB \leq_{\opjn} x$.
  \item pour tous $x,y\in X$, $x\opjn (x\opjn y) = (x\opjn x)\opjn y = x\opjn y$
    donc $x\opjn y$ est un majorant de $x$ (de même on prouve que $x\opjn y$
    est un majorant de $y$). Supposons que $z$ est un majorant de $\{x,y\}$ pour
    $\leq_{\opjn}$. Cela signifie que $x\opjn z = z$ et $y\opjn z = z$. Alors
    $(x\opjn y)\opjn z = x\opjn (y\opjn z) = x\opjn z = z$ donc
    $x\opjn y \leq_{\opjn} z$, donc $x\opjn y$ est inférieur à tout majorant de
    $\{x,y\}$~: c'est la borne supérieure de $\{x,y\}$.\qedhere
  \end{itemize}
\end{proof}

Ainsi on parlera de sup demi-treillis (respectivement inf demi-treillis) en
utilisant $\leq$ et les opérations $\opB,\opjn$ (respectivement $\opT,\opmt$)
sans souci, puisque toutes les définitions nous donnent l'ensemble de ces
notions.

\begin{exercise}
  Soit $\Sign$ une signature, on rappelle que l'algèbre de Lindenbaum-Tarski
  sur $\Sign$, notée $\LinTar(\Sign)$, est l'ensemble quotient
  $\quot{\Prop(\Sign)}{\dvvd}$ muni de l'ordre $\vdP$. Montrer que
  cette structure est un inf demi-treillis et un sup demi-treillis où le
  majorant est $\top$, le minorant $\bot$, la borne supérieure $\lor$ et la
  borne inférieure $\land$.
\end{exercise}

\begin{exercise}
  Soit un ensemble $X$, montrer que $(\powerset(X),\subseteq)$ est à la fois un
  sup et un inf demi-treillis, où $X$, $\vide$, $\cup$ et $\cap$ sont
  respectivement le majorant, le minorant, la borne supérieure et la borne
  inférieure.
\end{exercise}

On peut aussi définir les morphismes de demi-treillis, qui sont des morphismes
entre structures au sens attendu.

\begin{definition}[Morphisme de demi-treillis]
  Soient $(X,\leq)$ et $(Y,\leq)$ deux sup demi-treillis (respectivement deux
  inf demi-treillis). Un morphisme de sup demi-treillis (respectivement d'inf
  demi-treillis) est une fonction $\makeFunName{f}{X}{Y}$ croissante vérifiant
  l'égalité $f(\opB_X) = \opB_Y$ (respectivement $f(\opT_X) = \opT_Y$) et telle
  que, pour tous éléments $x,y\in X$, $f(x\opjn y) = f(x)\opjn f(y)$
  (respectivement $f(x\opmt y) = f(x)\opmt f(y)$).
\end{definition}

\subsection{Treillis}

Un demi-treillis ne concerne qu'une direction entre le sup et l'inf. Un
treillis, comme son nom l'indique, est une structure qui possède les propriétés
des deux demi-treillis à la fois.

\begin{definition}[Treillis]
  Un ensemble ordonné $(X,\leq)$ est un treillis si toute partie finie
  $F\subfin X$ possède à la fois une borne supérieure $\opJn F$ et une borne
  inférieure $\opMt F$.
\end{definition}

En théorie des ordres, il est direct de voir qu'être un treillis est équivalent
à être à la fois un inf demi-treillis et un sup demi-treillis, mais cela n'est
pas aussi évident algébriquement, car un inf demi-treillis est défini
algébriquement par la relation $\leq_{\opmt}$ et un sup demi-treillis par la
relation $\leq_{\opjn}$~: il faut vérifier que ces deux relations coïncident.

\begin{proposition}
  Un bimonoïde $(X,\opB,\opjn,\opT,\opmt)$ est un treillis pour $\leq_{\opjn}$
  (respectivement pour $\leq_{\opmt}$) si et seulement si les trois conditions
  suivantes sont vérifiées~:
  \begin{itemize}
  \item $(X,\opB,\opjn)$ est un inf demi-treillis,
  \item $(X,\opT,\opmt)$ est un inf demi-treillis,
  \item $\leq_{\opjn}$ et $\leq_{\opmt}$ coïncident
  \end{itemize}
  cette dernière condition étant équivalente aux deux identités d'absorption
  suivantes~:
  \begin{itemize}
    \begin{minipage}{0.4\textwidth}
    \item $\forall x,y\in X, (x\opjn y)\opmt x = x$
    \end{minipage}
    \begin{minipage}{0.4\textwidth}
    \item $\forall x,y\in X, (x\opmt y)\opjn x = x$
    \end{minipage}
  \end{itemize}
\end{proposition}

\begin{proof}
  Soit $(X,\opB,\opjn,\opT,\opmt)$ un bimonoïde commutatif idempotent. On
  suppose de plus que $\mathord{\leq_{\opjn}} = \mathord{\leq_{\opmt}}$,
  montrons les deux identités~:
  \begin{itemize}
  \item Soient $x,y\in X$, comme $x\leq_{\opjn} x\opjn y$, on en déduit que
    $x\leq_{\opmt} x\lor y$, c'est-à-dire que $x\opmt (x\opjn y) = x$ d'où
    l'égalité par commutativité de $\opmt$.
  \item Soient $x,y\in X$, comme $x\opmt y\leq_{\opmt} x$, on en déduit que
    $x\opmt y \leq_{\opjn} x$, c'est-à-dire que $(x\opmt y)\opjn x = x$.
  \end{itemize}

  Réciproquement, supposons vraies les identités. Montrons par double inclusion
  que $\mathord{\leq_{\opjn}} = \mathord{\leq_{\opmt}}$~:
  \begin{itemize}
  \item soient $x,y$ tels que $x\leq_{\opjn} y$, alors $x\opjn y = y$,
    montrons que $x\opmt y = x$~:
    \begin{align*}
      x \opmt y &= x \opmt (x \opjn y)\\
      &= x
    \end{align*}
    par la première identité. Ainsi $x\leq_{\opmt} y$, donc
    $\mathord{\leq_{\opjn}}\subseteq\mathord{\leq_{\opmt}}$.
  \item soient $x,y$ tels que $x\leq_{\opmt} y$, alors $x\opmt y = x$,
    montrons que $x\opjn y = y$~:
    \begin{align*}
      x \opjn y &= (x\opmt y) \opjn y\\
      &= y
    \end{align*}
    par la deuxième identité. Ainsi $x\leq_{\opjn} y$, donc
    $\mathord{\leq_{\opmt}}\subseteq\mathord{\leq_{\opjn}}$.\qedhere
  \end{itemize}
\end{proof}

Cela nous donne donc une définition purement algébrique d'un treillis (en
particulier, un treillis est un modèle d'une certaine théorie ne contenant que
l'égalité comme symbole de relation).

\begin{definition}[Morphisme de treillis]
  Soient $(X,\leq)$ et $(Y,\leq)$ deux treillis, un morphisme de treillis entre
  ces deux treillis est une fonction $\makeFunName{f}{X}{Y}$ qui est à la fois
  un morphisme de sup demi-treillis et un morphisme d'inf demi-treillis.
\end{definition}

\begin{exercise}
  Soit une signature $\Sign$, montrer que $\LinTar(\Sign)$ est un treillis.
\end{exercise}

\begin{exercise}
  Soit un ensemble $X$, montrer que $\powerset (X)$ est un treillis.
\end{exercise}

\begin{exercise}
  Soit une signature $\Sign$. Soit $\varMM$ une structure sur $\Sign$.
  Montrer que la fonction
  \[
  \makeFun{\Val}{\LinTar_1(\Sign)}{\powerset(\carMM)}
          {\clas{\varFF(x)}}{\compr{x \in \carMM}{\varMM\satisf\varFF(x)}}
  \]
  est un morphisme de treillis.
\end{exercise}

Cependant, la partie sup et la partie inf d'un treillis n'interagissent pas
suffisamment bien en général, c'est pourquoi il est nécessaire d'ajouter la
distributivité de $\opjn$ sur $\opmt$ et de $\opmt$ sur $\opjn$, si l'on veut
pouvoir l'utiliser.

\begin{definition}[Treillis distributif]
  Soit $(X,\leq)$ un treillis. On dit que c'est un treillis distributif si de
  plus l'une des deux propositions équivalentes est vraie~:
  \begin{enumerate}[label=(\roman*)]
  \item \label{enum.tr.distr1}
    $\forall x,y,z\in X, x\opmt (y\opjn z) = (x\opmt y)\opjn (x\opmt z)$
  \item \label{enum.tr.distr2}
    $\forall x,y,z\in X, x\opjn (y\opmt z) = (x\opjn y)\opmt (x\opjn z)$
  \end{enumerate}
\end{definition}

\begin{proof}
  Vérifions que ces deux propositions sont équivalentes. Par dualité, il nous
  suffit de prouver que $\ref{enum.tr.distr1} \implies\ref{enum.tr.distr2}$.
  On suppose donc \ref{enum.tr.distr1}. Soient $x,y,z\in X$, par un calcul~:
  \begin{align*}
    (x\opjn y) \opmt (x \opjn z) &= (x \opmt x) \opjn (y \opmt x)
    \opjn (x \opmt z) \opjn (y \opmt z)\\
    &= x\opjn (y\opmt x) \opjn (x \opmt z) \opjn (y \opmt z)\\
    &= x\opjn (y\opmt z) \opjn ((x\opmt y)\opjn (x\opmt z))\\
    &= x\opjn (y\opmt z) \opjn (x\opmt (y \opjn z))\\
    &= (x\opjn (x\opmt (y \opjn z)))\opjn (y \opmt z)\\
    (x\opjn y)\opmt (x\opjn z)&= x\opjn (y \opmt z)\qedhere
  \end{align*}
\end{proof}

\begin{exercise}
  Montrer que $\LinTar(\Sign)$ et $\powerset(X)$ sont des treillis
  distributifs (où $\Sign$ est une signature et $X$ un ensemble).
\end{exercise}

\subsection{Algèbre de Boole}

Nous voyons maintenant notre premier cas de structure logique décrite par un
ensemble ordonné. Avec la structure de treillis, nous avons naturellement une
structure permettant d'encoder la conjonction et la disjonction. Pour travailler
sur la logique propositionnelle, il ne nous manque que la négation,
l'implication pouvant se coder grâce à l'équivalence logique
$a\to b \equiv \lnot a \lor b$.

Nous souhaitons donc avoir, en plus de ce que nous avons déjà défini, une
fonction $\makeFunName{\opng}{X}{X}$ permettant de jouer le rôle de la négation.
En logique classique, qui est notre cadre de travail pour l'instant, deux
propriétés en particulier décrivent le comportement de $\lnot$~: le tiers exclu
et le principe de non contradiction, disant que $x\lor \lnot x$ est vraie et que
$x\land \lnot x$ est fausse. La notion de \og vrai\fg et de \og faux\fg sont
assez naturellement encodées par $\opT$ et $\opB$, respectivement, d'où la
définition suivante.

\begin{definition}[Complément]
  Soit un treillis $(X,\leq)$. On dit que $y$ est un complément de $x$ si
  \[x\opjn y = \opT \qquad x \opmt y = \opB\]
\end{definition}

\noindent
En fait, ceci suffit à caractériser la négation. En particulier, il n'y a qu'un
complément par élément, pour un treillis distributif.

\begin{proposition}
  Soit $(X,\leq)$ un treillis distributif et $x\in X$. Alors s'il existe un
  complément à $x$, celui-ci est unique.
\end{proposition}

\begin{proof}
  Supposons que $y,z$ soient deux compléments de $x$. Alors
  \begin{align*}
    z &= \opT \opmt z\\
    &= (x \opjn y) \opmt z\\
    &= (x \opmt z) \opjn (y \opmt z)\\
    &= \opB \opjn (y \opmt z)\\
    &= y \opmt z
  \end{align*}
  d'où $z\leq y$. En inversant le rôle de $y$ et $z$, on trouve aussi que
  $y\leq z$, d'où $y = z$.
\end{proof}

Une algèbre de Boole (nommée ainsi en l'honneur de George \name{Boole},
fondateur de l'approche algébrique de la logique \cite{boole54}) peut ainsi se
décrire grâce à cette notion de complément.

\begin{definition}[Algèbre de Boole]
  Une algèbre de Boole $(\varBB,\leq)$ est un treillis distributif
  complémenté. Pour un élément $x\in \varBB$, on notera $\opng x$ son
  complément.
\end{definition}

\begin{exercise}
  Soit une signature $\Sign$ et un ensemble $X$. Montrer que
  $\LinTar(\Sign)$ et $\powerset(X)$ sont des algèbres de Boole.
\end{exercise}

\begin{remark}
  En particulier, pour $X = \{\vide\}$, on obtient une algèbre de Boole
  de cardinal $2$, qui correspond en fait au corps $\quot{\bZ}{2\bZ}$.
  On notera cette algèbre de Boole $\varBB_2$, celle-ci a la
  particularité d'être la plus petite algèbre de Boole non dégénérée
  (telle que $\opB \neq \opT$) et d'être la seule algèbre de Boole qui a
  une structure de corps.
\end{remark}

\begin{exercise}
  Montrer que les identités de De Morgan sont vérifiées~:
  \[\opng \opT = \opB \qquad \opng\opng x = x
  \qquad \opng (x \opjn y) = \opng x \opmt \opng y
  \qquad \opng (x \opmt y) = \opng x \opjn \opng y\]
\end{exercise}

\begin{exercise}
  Montrer que si $\varBB$ et $\varBB'$ sont deux algèbres de Boole, alors
  un morphisme $f$ de treillis entre $\varBB$ et $\varBB'$ est en particulier un
  morphisme d'algèbres de Boole, au sens où la propriété suivante est
  vérifiée~:
  \[\forall x \in \varBB, f(\opng x) = \opng f(x)\]
\end{exercise}

L'exercice suivant donne une caractérisation purement algébrique d'une algèbre
de Boole~: on peut le voir comme un anneau idempotent.

\begin{exercise}
  Soit $(\varBB,\leq)$ une algèbre de Boole. On définit $\opXOR$ sur
  $\varBB$ par
  \[x\opXOR y \defeq (x\opjn y) \opmt (\opng x \opjn \opng y)\]
  Montrer que $(\varBB,\opB,\opT,\opXOR,\opmt)$ est un anneau (où $\opB$ est le
  neutre additif et $\opT$ le neutre multiplicatif).

  Réciproquement, soit $(\varAA,0,1,+,\times)$ un anneau idempotent,
  c'est-à-dire vérifiant
  \[\forall a\in \varAA, a\times a = a\]
  montrer les propositions suivantes~:
  \begin{enumerate}[label=(\roman*)]
  \item $\varAA$ est un anneau commutatif.
  \item $\varAA$ est un anneau de caractéristique $2$.
  \item la relation $\leq_{\opmt}$ définie par
    \[x\leq_{\opmt} y \defeq x\times y = x\]
    est une relation d'ordre.
  \item $(\varAA,\leq_{\opmt})$ est une algèbre de Boole, de majorant $1$, de
    minorant $0$, où $\times$ est la borne inférieure et où la borne supérieure
    est
    \[x\opjn y \defeq x + y + x \times y\]
  \end{enumerate}
\end{exercise}

\begin{remark}
  Chez certains auteurs, comme dans \cite{Halmos1966-HALLOB}, une algèbre de
  Boole est principalement traitée comme un anneau. Ce point de vue a
  l'avantage de simplifier les définitions, mais il permet moins de contrôle
  sur les hypothèses que l'on fait au niveau de la structure. Comme notre
  chapitre est dédié aux ensembles ordonnés, nous considérons plus intéressant
  de nous focaliser sur la structure ordonnée des algèbres de Boole pour
  en donner un traitement s'inscrivant dans une continuité, depuis les treillis
  jusqu'aux algèbres de Boole complètes.
\end{remark}

\subsection{Algèbre de Heyting}\label{subsec.heyting.alg}

Pour nous intéresser aux algèbres de Heyting, il nous faut d'abord parler
de la logique intuitionniste. Pour ne pas nous attarder sur les questions
philosophiques qui sont à l'origine de l'essor des mathématiques intuitionnistes
et constructives, nous nous contenterons de remarquer que c'est une logique plus
faible, mais donnant de meilleurs résultats (les résultats sont plus durs à
montrer mais apportent plus d'information).

La logique intuitionniste naît d'une volonté de travailler sans le tiers exclu.
D'un point de vue algorithmique, par exemple, ce principe peut être remis en
question~: comment pouvons-nous effectivement décider toute proposition~? Il
est en effet difficile d'imaginer une machinerie finie qui pourrait décider,
par exemple, si tel objet est infini. Ainsi une part des mathématiques n'a ni
besoin, ni envie du principe du tiers exclu, qui est équivalent au raisonnement
par l'absurde. Nous étudierons en détail cette logique lors de notre étude de la
théorie de la démonstration (notamment aux \cref{chp.sequents} et
\cref{chp.deduc.nat}), mais nous donnons ici une définition élémentaire de la
logique intuitionniste.

\begin{definition}
  Soit une signature $\Sign$. On définit la relation de déduction syntaxique
  $\mathord{\vdNJ}\subseteq\Ctx(\Sign)\times\Formula(\Sign)$
  par les mêmes règles que pour $\vdP$, à l'exception de $\RbotC$ que l'on
  remplace par la règle suivante~:
  \begin{prooftree}
    \AxiomC{$\varCG\vdNJ \bot$}
    \RightLabel{$\RbotE$}
    \UnaryInfC{$\varCG\vdNJ \varFF$}
  \end{prooftree}
  appelée communément principe d'explosion, ou bien par sa dénomination latine
  \latinexpr{ex falso quodlibet} (version raccourcie de
  \latinexpr{ex falso sequitur quodlibet}).
\end{definition}

\begin{exercise}\label{exo.EDN.RAA}
  Soit la règle
  \begin{center}
    \AxiomC{}
    \RightLabel{$\RLEM$}
    \UnaryInfC{$\varCG\vdP \varFF\lor \lnot \varFF$}
    \DisplayProof
  \end{center}
  Montrer que celle-ci est équivalente à $\RbotC$ modulo $\vdNJ$,
  c'est-à-dire que si l'on considère $\vdNJ$ et qu'on rajoute l'une des deux
  règles, alors la seconde est dérivable.
\end{exercise}

\begin{exercise}\label{exo.LinTar.LJ}
  Vérifier que dans le cas intuitionniste aussi, la relation $\vdNJ$ est une
  relation de pré-ordre. On notera $\LinTarInt(\Sign)$ l'algèbre de
  Lindenbaum-Tarski intuitionniste associée à $\Sign$.
\end{exercise}

\begin{exercise}
  Montrer que $\LinTarInt(\Sign)$ est un treillis distributif.
\end{exercise}

On peut alors se demander quelle est la différence entre ce que nous pouvons
prouver en logique intuitionniste et ce que nous pouvons prouver en logique
classique. La réponse est difficile à donner directement, mais un point
important est que ce que l'on peut prouver en logique intuitionniste est
constructif~: par exemple, prouver qu'un objet existe signifie qu'on a
effectivement exhibé un objet. En comparaison, on peut prouver grâce à l'absurde
qu'un objet existe en prouvant simplement qu'il ne peut pas ne pas exister.

Une conséquence de cet effet est que certaines tautologies ne sont plus valides,
en particulier la tautologie
\[\lnot(\forall x, \lnot \varFF) \implies \exists x, \varFF\]
exprimant qu'à partir de l'impossibilité que tous les $x$ rejettent $\varFF$,
on peut extraire un certain $x$ acceptant $\varFF$.

De même pour la tautologie
\[\lnot (\varFF \land \varFP) \implies \lnot \varFF \lor \lnot \varFP\]
car pour la prouver, il faudrait de façon uniforme décider si $\lnot \varFF$ est
vraie ou si $\lnot \varFP$ est vraie. Comme cela dépend directement que la
valeur de vérité de $\varFF$ et de celle de $\varFP$, il est impossible de
conclure en logique intuitionniste. En échange, si l'on prouve
$\vdNJ \varFF \lor \varFP$, cela signifie qu'on
peut soit prouver $\vdNJ \varFF$, soit prouver $\vdNJ \varFP$
(\latinexpr{c.f.} \cref{prop.disjonction.NJ}).

La question, maintenant, est de savoir s'il existe une structure similaire aux
algèbre de Boole, mais dans laquelle la logique appliquée est la logique
intuitionniste. C'est le cas des algèbres de Heyting, dont la définition repose
avant tout sur celle de l'opération $\opimp$, d'où l'on dérive le sens de
$\opng$ par $\opng x = x \opimp \opB$, comme nous l'avons déjà fait.

\begin{definition}[Pseudo-complément relatif]
  Soit un treillis distributif $(X,\leq)$, et deux éléments $a,b\in X$. On dit
  que $x$ est le pseudo-complément de $a$ relatif à $b$ si $x$ est le
  maximum de l'ensemble $\compr{x \in X}{a \opmt x \leq b}$. Si le
  pseudo-complément de $a$ relatif à $b$ existe, on le note $a\opimp b$.
  On appelle pseudo-complément de $a$ le pseudo complément de $a$ relatif à
  $\opB$, s'il existe, et on le note alors $\opng x$.
\end{definition}

\begin{definition}[Algèbre de Heyting \cite{heyting30}]
  Un treillis distributif $(\varBH,\leq)$ est une algèbre de Heyting si
  tout élément admet un pseudo-complément relatif à tout autre élément,
  c'est-à-dire si pour tous $a,b\in \varBH$, l'ensemble
  $\compr{x\in \varBH}{a \opmt x \leq b}$ admet un
  maximum.
\end{definition}

\begin{remark}
  Une inégalité par rapport à un pseudo-complément peut se réécrire par une
  conjonction~: $a\leq b \opimp c$ revient directement à $a\opmt b \leq c$ étant
  donné que $b \to c$ est le maximum des éléments $x$ tels que
  $x\opmt b \leq c$. On peut donc formaliser le modus ponens~: si l'on a
  $a\leq b \opimp c$ et $a\leq b$, alors $a\opmt b \leq c$ et $a\opmt b = a$,
  donc $a\leq c$.
\end{remark}

\begin{exercise}
  Soit une signature $\Sign$. Montrer que $\LinTarInt(\Sign)$ est
  bien une algèbre de Heyting.
\end{exercise}

\begin{exercise}
  Montrer qu'une algèbre de Boole est en particulier une algèbre de Heyting, et
  que le complément d'un élément est en particulier son pseudo-complément.
\end{exercise}

\subsection{Treillis complet}

Notre première définition d'un treillis peut paraître légèrement plus
artificelle que sa caractérisation algébrique. En effet, dans la version
algébrique, on peut voir un treillis comme un certain bimonoïde, que l'on peut
ensuite étendre en l'existence d'une borne supérieure et d'une borne inférieure
pour toute partie finie (en raisonnant juste par récurrence). A l'inverse,
lorsque l'on part de la définition de treillis comme ensemble ordonné possédant
toutes les bornes sur des parties finies, une question naît immédiatement~:
pourquoi ne pas autoriser des parties infinies~? C'est pour intégrer cette
possibilité que nous introduisons la notion de treillis complet.

\begin{definition}[Treillis complet]
  Un treillis complet est un ensemble ordonné $(X,\leq)$ tel que toute partie
  $Y\subseteq X$ possède à la fois une borne supérieure et une borne inférieure.
\end{definition}

Un treillis complet est, par définition, plus fort qu'un treillis. Dans la
pratique, on trouve effectivement plusieurs treillis qui ne sont pas complets
(le fait d'avoir deux notions n'est donc pas une distinction simplement formelle
mais revêt une vraie différence conceptuelle). Remarquons aussi qu'un treillis
complet peut se caractériser par l'existence d'une seule borne.

\begin{property}\label{prop.treillis.complet.justInf}
  L'ensemble ordonné $(X,\leq)$ est un treillis complet si et seulement si toute
  partie possède une borne supérieure.
\end{property}

\begin{proof}
  Un sens est évident. Pour l'autre, supposons que toute partie de $X$ possède
  une borne supérieure, et montrons alors qu'elle possède une borne inférieure.

  Soit $\varPt\subseteq X$, on définit $\downcl{\varPt}$ par
  \[\downcl{\varPt} \defeq \compr{x\in X}{\forall y \in \varPt, x \leq y}\]
  et soit $y_0$ la borne supérieure de $\downcl{\varPt}$. Montrons que $y_0$ est
  la borne inférieure de $\varPt$. Tout d'abord, si $y\in \varPt$, alors par
  définition, $y$ est un majorant de $\downcl{\varPt}$, donc puisque $y_0$ est
  le plus petit majorant, $y_0$ est un minorant de $\varPt$. De plus, si $y$ est
  un minorant de $\varPt$, cela signifie que $y\in \downcl{\varPt}$, donc comme
  $y_0$ est un majorant de $\downcl{\varPt}$, on en déduit que $y\leq y_0$.
  Ainsi $y_0$ est le plus grand minorant de $\varPt$~: c'est la borne inférieure
  de $\varPt$.
\end{proof}

On peut maintenant prouver l'un des théorèmes importants que nous avons présenté
dans le \cref{chp.induction}~: le \cref{thm.ind.KT}, mais dans une version bien
plus forte et générale.

\begin{theorem}[Knaster-Tarski \cite{KTFull}]\label{thm.KT}
  Soit $(X,\leq)$ un treillis complet et $\makeFunName{f}{X}{X}$ une fonction
  croissante. Alors l'ensemble
  \[\fix(f) \defeq \compr{x\in X}{f(x) = x}\]
  est un treillis complet pour l'ordre induit par $\leq$.
\end{theorem}

\begin{proof}
  On définit tout d'abord l'ensemble des pré points fixes de $f$~:
  \[\prefix(f) \defeq \compr{x\in X}{f(x)\leq x}\]
  et on montre qu'il est un treillis complet pour $\leq$. Soit une partie
  $\varPt\subseteq\prefix(f)$, montrons que $\varPt$ possède une borne
  inférieure. Nous prenons pour candidat de la borne inférieure de $\varPt$ dans
  $\prefix(f)$, la borne inférieure $\opMt \varPt$ dans $X$~: montrons donc que
  $\opMt \varPt$ est bien un pré point fixe de $f$ (cela suffira à montrer qu'on
  a une borne inférieure, puisque la borne inférieure dans $\varPt$ est aussi la
  borne inférieure dans $\prefix(f)$, à partir du moment où la borne appartient
  effectivement à $\prefix(f)$).

  On montre que $f\left(\opMt \varPt\right) \leq \opMt \varPt$. Pour cela, il
  nous suffit de montrer que $f\left(\opMt \varPt\right)$ est un minorant de
  $\varPt$, puisque l'on sait que $\opMt \varPt$ est le plus grand minorant.
  Soit $y\in \varPt$, on remarque que $\opMt \varPt \leq y$, donc
  $f\left(\opMt \varPt\right) \leq f(y)$, mais comme par définition $f(y)\leq y$
  et par transitivité, on en déduit que $f\left(\opMt \varPt\right) \leq y$.
  Ainsi $f\left(\opMt \varPt\right) \leq \opMt \varPt$.

  On sait donc que toute partie de $\prefix(f)$ admet une borne inférieure. Par
  la propriété précédente (son dual, plus précisément), $\prefix(f)$ est donc
  un treillis complet pour $\leq$. On a donc un treillis complet
  $(\prefix(f),\leq)$ et une fonction $\makeFunName{f}{\prefix(f)}{\prefix(f)}$.
  En prenant le dual, on obtient donc que $\postfix(\prefix(f))$ est un treillis
  complet.

  Montrons maintenant que $\postfix(\prefix(f)) = \fix(f)$. Par définition,
  \[\postfix(\prefix(f)) =
  \compr{x\in \compr{x\in X}{f(x)\leq x}}{x \leq f(x)}\]
  ce qui correspond exactement à l'ensemble
  $\compr{x\in X}{(f(x) \leq x) \land (x \leq f(x))}$. Comme $\leq$ est
  antisymétrique, cet ensemble est $\fix(f)$. Ainsi $\fix(f)$ est un treillis
  complet.
\end{proof}

\begin{remark}
  On en déduit donc, en particulier, qu'il existe un plus petit et un plus grand
  point fixe pour $f$. En rentrant un peu plus en détail dans la démonstration,
  on peut vérifier que le plus petit point fixe est le plus petit pré point
  fixe, et que le plus grand point fixe est le plus grand post point fixe.
\end{remark}

\begin{exercise}
  Montrer que $(\powerset(X),\subseteq)$ est un treillis complet, pour tout
  ensemble $X$. En déduire le \cref{thm.ind.KT}.
\end{exercise}

Enfin, donnons deux mots à propos des algèbre de Heyting complètes.
D'abord, précisons la définition de cette structure~:

\begin{definition}[Algèbre de Heyting complète]
  Un ensemble ordonné $(\varBH,\leq)$ est une algèbre de Heyting complète
  si $(\varBH,\leq)$ est un treillis complet et si toute paire
  $(a,b)\in \varBH^2$ possède un pseudo-complément relatif.
\end{definition}

Un point important pour ces structure est que l'existence des pseudo-compléments
apporte une nouvelle loi de distributivité, infinie. En effet, étant donnée une
algèbre de Heyting complète $\varBH$, une partie $\varPt\subseteq \varBH$ et un
élément $y\in \varBH$, on a l'identité~:
\[\left(\opJn_{x\in \varPt} x\right)\opmt y = \opJn_{x \in \varPt}(x \opmt y)\]
Cette identité permet de récupérer directement la distributivité, et fonctionne
en fait aussi dans le cas d'une algèbre de Heyting non complète (sous
réserve d'existence des bornes supérieures), ce qui fait qu'on peut affaiblir
les définitions d'algèbre de Heyting en un simple treillis contenant un
pseudo-complément relatif. On peut trouver des détails sur les liens entre ces
propriétés dans \cite{lattice}.

\section{Filtres}

Enfin, présentons la notion de filtre, et la notion duale d'idéal. Nous allons
détailler comment se comportent les filtres dans les structure de plus en plus
fortes~: demi-treillis, treillis et algèbre de Boole, puis nous verrons
les ultrafiltres et le lemme de l'ultrafiltre.

Les filtres ont été introduits à l'origine par Henri \name{Cartan} dans
\cite{bourbaki1971topologie}, dans un but topologique (nous verrons dans le
\cref{chp.topo} que beaucoup d'outils topologiques peuvent s'interpréter à
travers la notion de filtre). Ces filtres étaient alors définis sur les parties
d'un ensemble, mais la présentation que nous en faisons ici est plus générale
puisqu'elle part d'un ensemble ordonné quelconque, en cherchant à ajouter avec
parcimonie des hypothèses, jusqu'à arriver à des filtres sur des algèbres de
Boole.

\subsection{Définitions et caractérisations}

\noindent
Commençons par définir ce qu'est un filtre.

\begin{definition}[Filtre]
  Soit $(X,\leq)$ un ensemble ordonné. Un filtre est une partie
  $\varFil\subseteq X$ vérifiant~:
  \begin{itemize}
  \item $\varFil$ est non vide~: $\varFil \neq\vide$.
  \item $\varFil$ est clos par le haut~: pour tout $x\in \varFil$ et
    $y\in X$, si $x\leq y$ alors $y\in \varFil$.
  \item tous les éléments de $\varFil$ sont compatibles vers le bas~: pour
    tous $x,y\in\varFil$, il existe $z\in\varFil$ tel que $z\leq x$ et
    $z\leq y$.
  \end{itemize}
  Un filtre $\varFil$ est dit propre si $\varFil \neq X$.
\end{definition}

Un filtre peut donc être vu comme une description d'une notion de grandeur~:
un élément d'un filtre est vu comme un \og grand\fg élément. En plus de cela,
cette notion de grandeur doit avoir une notion de compatibilité~: deux grands
éléments doivent avoir en commun un grand élément. La notion duale, qui peut
donc être vue comme la spécification d'une notion de petitesse, est celle
d'idéal.

\begin{definition}[Idéal]
  Soit $(X,\leq)$ un ensemble ordonné. Un idéal est une partie
  $\varId \subseteq X$ vérifiant~:
  \begin{itemize}
  \item $\varId$ est non vide~: $\varId \neq\vide$.
  \item $\varId$ est clos par le bas~: pour tout $x\in \varId$ et
    $y\in X$, si $y\leq x$ alors $y\in \varId$.
  \item tous les éléments de $\varId$ sont compatibles vers le haut~:
    pour tous $x,y\in \varId$, il existe $z\in \varId$ tel que
    $x\leq z$ et $y\leq z$.
  \end{itemize}
  Un idéal $\varId$ est dit propre si $\varId \neq X$.
\end{definition}

\begin{exercise}
  Soit une algèbre de Boole $(\varBB,\leq)$. Montrer qu'une partie
  $\varId \subseteq \varBB$ est un idéal de $\varBB$ si et seulement si c'est
  un idéal au sens algébrique de l'anneau associé
  $(\varBB,\opB,\opT,\opXOR,\opmt)$.
\end{exercise}

Nous ne travaillerons par sur les idéaux, mais donnons la définition car
celle-ci reste importante.

\begin{example}
  Soit un ensemble $X$ et un élément $x\in X$. On peut définir le filtre
  principal sur $\powerset(X)$ en $x$, $\FilPrinc{x}$, par
  \[\FilPrinc{x} \defeq \compr{\varPt\subseteq X}{x \in \varPt}\]

  On voit que $X\in \FilPrinc{x}$, que si $\varPt\subseteq \varPtB$ et
  $x\in \varPt$ alors $x\in \varPtB$, et que si $x\in \varPt$ et $x\in \varPtB$
  alors $x\in \varPt\cap \varPtB$. Ainsi l'ensemble ordonné $\powerset(X)$
  contient au moins autant de filtres que d'éléments dans $X$.

  Pour une partie $\varPt \subseteq X$, on peut aussi définir
  \[\FilPrinc{\varPt} \defeq \compr{\varPtB}{\varPt\subseteq\varPtB}\]
  dont on vérifie de la même manière qu'il est un filtre. Le filtre principal en
  $x$ est ainsi $\FilPrinc{\{x\}}$.
\end{example}

On peut donner une première caractérisation des filtres, dans le cas d'un
inf demi-treillis.

\begin{property}
  Soit $(X,\leq)$ un inf demi-treillis de majorant $\opT$ et de borne inf
  $\opmt$. Alors $\varFil\subseteq X$ est un filtre sur $X$ si et seulement
  si les conditions suivantes sont vérifiées~:
  \begin{itemize}
    \begin{minipage}{0.2\textwidth}
    \item $\opT \in X$~;
    \end{minipage}
    \begin{minipage}{0.35\textwidth}
    \item $\varFil$ est clos par le haut~;
    \end{minipage}
    \begin{minipage}{0.3\textwidth}
    \item si $x,y\in \varFil$, alors $x\opmt y\in \varFil$.
    \end{minipage}
  \end{itemize}
\end{property}

\begin{proof}
  Supposons que $\varFil$ vérifie les propriétés listées. Le fait d'être
  clos par le haut est par définition. Comme $\opT\in\varFil$,
  $\varFil$ est non vide. Si $x,y\in\varFil$, alors il existe
  effectivement $z\in \varFil$ tel que $z\leq x$ et $z\leq y$, en prenant
  $z = x\opmt y$.

  Réciproquement, supposons que $\varFil$ est un filtre. Alors, puisqu'il
  possède un élément, disons $x$, et que $x\leq \opT$, on en déduit par clôture
  par le haut que $\opT\in\varFil$. Le fait que $\varFil$ est clos par le
  haut est par hypothèse. Soient $x,y\in \varFil$, alors par hypothèse il
  existe $z\in \varFil$ tel que $z\leq x$ et $z\leq y$. Par définition de la
  borne inférieure, on a donc que $z\leq x\opmt y$, et par clôture par le haut
  on en déduit que $x\opmt y \in \varFil$.
\end{proof}

Les inf demi-treillis permettent donc de mieux caractériser les filtres. En
quelque sorte, puisqu'ils permettent de délimiter le fait de devenir plus petit,
il est plus aisé d'y définir ce qui est grand (c'est-à-dire ce qui arrive à
rester grand quand on le rend plus petit). Les sup demi-treillis, eux, ne
permettent pas de mieux définir un treillis, mais de mieux en délimiter les
caractéristiques.

\begin{property}
  Soit $(X,\leq)$ un sup demi-treillis et $\varFil \subseteq X$ un filtre.
  Alors $\varFil$ est propre si et seulement si $\opB\notin\varFil$.
\end{property}

\begin{proof}
  Puisque $\varFil$ est clos par le haut, si $\opB\in\varFil$ alors tout
  élément est dans $\varFil$, donc $\varFil$ n'est pas propre (et si
  $\varFil$ n'est pas propre, alors il contient $\opB$).
\end{proof}

\noindent
Les sup demi-treillis permettent de définir la primalité, propriété d'intérêt
pour les filtres~:

\begin{definition}[Filtre premier]
  Soit $(X,\leq)$ un sup demi-treillis et $\varFil$ un filtre sur $X$. On dit
  que $\varFil$ est premier s'il est propre et si pour tous $x,y\in X$, si
  $x\opjn y \in \varFil$ alors $x\in F$ ou $y\in F$.
\end{definition}

On peut voir une analogie ici avec un idéal dans un anneau~: un idéal premier
est un idéal tel que, lorsqu'on a un produit $a\times b$ dedans, celui-ci doit
nécessairement provenir du produit par un élément de l'idéal. De la même façon,
ici, on sait naturellement que si $x\in\varFil$ alors
$x\opjn y \in\varFil$, et la notion de primalité signifie que tout élément de
la forme $x\opjn y$ doit provenir d'un cas de ce type, où $x$ (ou $y$) était
déjà dans $\varFil$.

Pour synthétiser les différentes propriétés dans le cas d'un treillis~: les
filtres propres (qui sont les filtres intéressants) sont les parties closes par
le haut qui possèdent $\opT$ mais pas $\opB$, et qui sont closes par
intersection finie.

Supposons maintenant qu'on possède une partie $X$ (non vide) qui n'est pas un
filtre. Comment trouver un filtre correspondant naturellement à $X$ ?
Premièrement, il nous faut d'abord clore $X$ par le haut~: on prend l'ensemble
des éléments supérieurs à au moins un élément de $X$. Il nous faut aussi nous
assurer que pour tout $x,y$ dans notre filtre, un élément inférieur à $x$ et $y$
soit dans le filtre. Le plus simple ici est de considérer un inf demi-treillis~:
on veut ajouter tous les $x\opmt y$, mais comme il est possible d'itérer ces
bornes inférieures, il nous faut ajouter tous les $\opMt F$ pour
$F\subfin X$. On obtient ainsi la construction suivante~:

\begin{definition}[Filtre engendré par une partie]
  Soit $(X,\leq)$ un inf demi-treillis et $\varPt\subseteq X$ une partie de $X$.
  On définit
  \[\filclose{\varPt} \defeq
  \compr{x\in X}{\exists F\subfin \varPt, \opMt F \leq x}\]
\end{definition}

\begin{proposition}
  Soit $(X,\leq)$ un inf demi-treillis et $\varPt\subseteq X$ une partie de $X$.
  Alors $\filclose{\varPt}$ est un filtre sur $X$.
\end{proposition}

\begin{proof}
  Tout d'abord, on remarque que $\opT = \opMt \vide$ et
  $\vide\subfin \varPt$, donc $\filclose \varPt$ est non vide. Soit
  $x\in \filclose{\varPt}$ et $y\in X$ tel que $x\leq y$. On trouve
  $F\subfin \varPt$ tel que $\opMt F\leq x$, alors $\opMt F \leq y$ par
  transitivité, donc $y\in \filclose \varPt$. Soient $x$ et $y$ des éléments de
  $\filclose \varPt$, on trouve $F$ et $F'$ des parties finies de $\varPt$
  telles que $\opMt F \leq x$ et $\opMt F' \leq y$. Alors
  \[\opMt (F\cup F') = \left(\opMt F\right)\land\left(\opMt F\right)
  \leq x\opmt y\]
  donc $x\opmt y \in \filclose{\varPt}$. Ainsi $\filclose{\varPt}$ est un
  filtre.
\end{proof}

\begin{remark}
  Dans le cas où $\varPt = \vide$, $\filclose \varPt = \{\opT\}$.
\end{remark}

Par défaut, le filtre engendré par une partie peut tout à fait contenir tous
les éléments. Prenons par exemple une algèbre de Boole $X$, un élément $x$ et le
filtre engendré par $\{x,\opng x\}$~: il est évident que ce filtre n'est pas
propre. Cependant, il existe une propriété assurant qu'un filtre engendré est
propre, et qui est tout à fait naturelle~: on veut que les $\opMt F$ soient
tous différents de l'élément minimum. \'Evidemment, cette propriété n'a
de sens que dans le cas d'un treillis.

\begin{definition}[Propriété de l'intersection finie]
  Soit $(X,\leq)$ un treillis et $\varPt\subseteq X$. On dit que $\varPt$ a la
  propriété de l'intersection finie si
  \[\forall F\subfin \varPt, \opMt F\neq \opB\]
\end{definition}

\begin{proposition}
  Soit $(X,\leq)$ un treillis et $\varPt\subseteq X$ une partie non vide de $X$.
  Alors $\filclose \varPt$ est un filtre propre si et seulement si $\varPt$ a la
  propriété de l'intersection finie.
\end{proposition}

\begin{proof}
  Par définition, $\opB\in \filclose{\varPt}$ si et seulement s'il existe
  $F\subfin \varPt$ tel que $\opMt F \leq \bot$, mais la condition que
  $\opMt F\leq \opB$ est équivalente à ce que $\opMt F = \bot$, étant donné que
  l'autre inégalité est toujours vérifiée. Ainsi $\opB\in\filclose{\varPt}$ si
  et seulement s'il existe $F\subfin \varPt$ tel que $\opMt F = \opB$, donc
  $\filclose{\varPt}$ est propre si et seulement si $\varPt$ possède la
  propriété de l'intersection finie.
\end{proof}

\begin{exercise}
  Soit $\Sign$ une signature. On considère $\varT\in\Theory(\Sign)$
  une théorie. Montrer que $\varT$ est cohérente (c'est-à-dire que
  $\varT\nvdP \bot$) si et seulement si $\varT'$, la projection
  de $\varT$ dans $\LinTar(\Sign)$, a la propriété de l'intersection
  finie.
\end{exercise}

\begin{exercise}
  Soit $\Sign$ une signature. Soit $\varT\in\Theory(\Sign)$ une
  théorie. Montrer que la projection de $\vdTclos$ dans
  $\LinTar(\Sign)$ est le filtre engendré par la projection de $\varT$
  dans $\LinTar(\Sign)$.
\end{exercise}

\begin{exercise}
  On se place maintenant sur un ensemble $X$ infini quelconque, et on considère
  le filtre de Fréchet\footnote{Nommé en l'honneur de Maurice \name{Fréchet}}
  sur $X$~:
  \[\Frechet \defeq \compr{\varPt\subseteq X}
             {X \setminus \varPt\text{ est fini}}\]

  Montrer que le filtre de Fréchet est bien un filtre. Montrer que pour
  tout $x \in X$, $\FilPrinc{x}$ ne contient pas $\Frechet$.
\end{exercise}

\subsection{Ultrafiltre}\label{sbsct.ultrafiltre}

On définit maintenant la notion d'ultrafiltre, qui est la notion duale de celle
d'idéal maximal, souvent rencontrée en algèbre.

\begin{definition}[Ultrafiltre]
  Soit $(X,\leq)$ un ensemble ordonné. On dit qu'un filtre $\varFil$ est un
  ultrafiltre s'il est un filtre propre maximal pour l'inclusion, c'est-à-dire
  si pour tout autre filtre propre $\varFilB$ de $X$, si
  $\varFil \subseteq\varFilB$ alors $\varFil = \varFilB$.
\end{definition}

\begin{remark}
  L'inclusion entre filtres est aussi appelée \og finesse\fg~: un filtre
  $\varFil$ est dit plus fin qu'un filtre $\varFilB$ si
  $\varFil\supseteq\varFilB$. L'idée derrière cette appellation est qu'un
  filtre plus gros va définir une notion de grandeur plus petite (puisque les
  éléments que l'on peut ajouter à un filtre seront par définition plus petits).
  Un ultrafiltre est donc un filtre propre le plus fin possible.
\end{remark}

\noindent
Dans le cas d'un treillis distributif, un ultrafiltre est en particulier
premier.

\begin{proposition}
  Soit $(X,\leq)$ un treillis distributif. Si un filtre $\varFil$ est un
  ultrafiltre, alors c'est un filtre premier.
\end{proposition}

\begin{proof}
  Supposons que $\varFil$ est un ultrafiltre, montrons que c'est alors un
  filtre premier. Supposons par l'absurde qu'il existe $x,y\in X$ tels que
  $x\opjn y \in \varFil$, $x\notin\varFil$ et $y\notin\varFil$.
  Montrons alors que soit $\varFil\cup\{x\}$ soit $\varFil\cup\{y\}$ a la
  propriété de l'intersection finie.

  Supposons qu'on trouve $F\subfin \varFil$ tel que
  $\opMt F \opmt x =\opB$ et qu'on trouve $F'\subfin\varFil$ tel que
  $\opMt F' \opmt y = \opB$. En prenant $G \defeq F \cup F'$, on a donc
  \[\opMt G \land x = \opB \qquad \opMt G \opmt y = \opB\]
  donc $\left(\opMt G \opmt x\right)\opjn\left(\opMt G\opmt y\right) = \opB$,
  et en développant par distributivité on obtient donc
  \[\left(\opMt G \opjn y\right)\opmt\left(\opMt G\opjn \opMt G\right)
  \opmt \left(x\opjn \opMt G\right)\opmt(x\opjn y) = \opB\]
  comme $\opMt G\in \varFil$, on en déduit que
  $\opMt G \opjn y \in \varFil$ et $\opMt G \opjn x \in \varFil$,
  et par hypothèse $x\opjn y \in \varFil$. Ainsi, par stabilité de
  $\varFil$ par $\opmt$, et grâce à l'égalité, on en déduit que
  $\opB\in\varFil$, ce qui contredit l'hypothèse que $\varFil$ est un
  ultrafiltre (donc propre).

  On en déduit, sans perte de généralité, que $\varFil\cup\{x\}$ a la
  propriété de l'intersection finie, ce qui signifie que
  $\filclose{\varFil\cup\{x\}}$ est un filtre, contenant par définition
  $\varFil$~: comme $\varFil$ est supposé maximal, on en déduit que
  $x\in \varFil$. Ainsi $\varFil$ est premier.
\end{proof}

\noindent
On a, de plus, une caractérisation encore meilleure dans le cas d'une algèbre de
Boole.

\begin{proposition}
  Soit $(\varBB,\leq)$ une algèbre de Boole et $\varFil$ un filtre. Les
  conditions suivantes sont équivalentes~:
  \begin{enumerate}[label=(\roman*)]
  \item \label{enum.ultra.prim1} $\varFil$ est un ultrafiltre.
  \item \label{enum.ultra.prim2} $\varFil$ est premier.
  \item \label{enum.ultra.prim3} pour tout $x\in \varBB$,
    $x\in \varFil \iff \opng x \notin\varFil$.
  \end{enumerate}
\end{proposition}

\begin{proof}
  On va montrer le cycle d'implications
  $\ref{enum.ultra.prim1} \implies\ref{enum.ultra.prim2}
  \implies\ref{enum.ultra.prim3}\implies\ref{enum.ultra.prim1}$
  dont on a déjà prouvé la première implication. Il nous reste à montrer les
  deux autres.
  
  Si $\varFil$ est un filtre premier, alors comme $\opT \in \varFil$ et
  $\opT = x \opjn \opng x$, on sait que soit $x$ soit $\opng x$ est dans
  $\varFil$. Comme par définition $\opB\notin\varFil$ et que
  $x\opmt\opng x = \opB$, on en déduit qu'un seul des deux est dans
  $\varFil$, ce qui est le résultat qu'on veut démontrer.

  Si pour tout $x\in \varBB$, $x\in\varFil\iff\opng x\notin\varFil$ et
  si $\varFilB$ est tel que $\varFil\subseteq\varFilB$ et $\varFilB$
  est un filtre propre, alors si l'inclusion est stricte, on trouve $x$ un
  élément de $\varFilB$ qui n'est pas dans $\varFil$. Mais par définition,
  puisque $x\notin\varFil$, $\opng x\in\varFil$, or $x\in \varFilB$ et
  par inclusion $\opng x \in \varFilB$, donc $\opB\in \varFilB$, ce qui
  contredit le fait que $\varFilB$ est propre. On en déduit que
  $\varFilB = \varFil$, et donc que $\varFil$ est un ultrafiltre.
\end{proof}

\begin{exercise}
  Soit $X$ un ensemble et $x\in X$. Montrer que $\FilPrinc{x}$ est un
  ultrafiltre dans l'algèbre de Boole $(\powerset(X),\subseteq)$.
\end{exercise}

Pour finir, nous allons démontrer le lemme de l'ultrafiltre, qui est une
conséquence du lemme de Zorn. Il dit que tout filtre propre peut être
étendu en un ultrafiltre.

\begin{theorem}[Lemme de l'ultrafiltre]\label{thm.ultrafilter.lemma}
  Soit $(X,\leq)$ un treillis. Soit $\varFil$ un filtre propre sur $X$. Alors il
  existe un ultrafiltre $\varUlF$ contenant $\varFil$.
\end{theorem}

\begin{proof}
  On définit l'ensemble ordonné
  \[\varXX \defeq \compr{\varFilB \subseteq X}{\varFilB
  \text{ est un filtre propre}}\]
  muni de l'inclusion ensembliste.

  On veut appliquer le lemme de Zorn sur cet ensemble. Pour cela, on
  vérifie que $\varXX$ est inductif.

  Soit $\{\varFil_i\}_{i\in I}$ une chaîne d'éléments de $\varXX$,
  montrons que $\varFilB \defeq \displaystyle\bigcup_{i\in I} \varFil_i$ est un
  majorant de cette chaîne (si $I \neq \vide$, car sinon n'importe quel
  élément est un majorant de la chaîne). Il nous suffit pour cela de vérifier
  que $\varFilB$ est bien un filtre propre~:
  \begin{itemize}
  \item soit $i\in I$, comme $\opT \in \mathcal F_i$, on en déduit que
    $\opT \in \varFilB$.
  \item soit $x\in \varFilB$ et $y\in X$ tel que $x\leq y$. Par définition
    de $\varFilB$, on trouve $i\in I$ tel que $x\in \varFil_i$, et comme
    $\varFil_i$ est un filtre on en déduit que $y\in \varFil_i$, d'où
    $y\in\varFilB$.
  \item soient $x,y\in \varFilB$. On trouve $i,j\in I$ tels que
    $x\in \varFil_i$ et $y\in \varFil_j$. Comme $\{\varFil_i\}_{i \in I}$ est une
    chaîne pour l'inclusion, alors soit $\varFil_i \subseteq \varFil_j$,
    soit $\varFil_j\subseteq \varFil_i$. Sans perte de généralité, on
    suppose donc que $\varFil_i\subseteq \varFil_j$, donc que
    $x\in \varFil_j$. Comme $\varFil_j$ est un filtre, on en déduit que
    $x\opmt y \in \varFil_j$, donc que $x\opmt y \in \varFilB$.
  \item supposons que $\opB\in \varFilB$. Alors on trouve $i\in I$ tel que
    $\opB\in \varFil_i$, mais comme tous les $\varFil_i$ sont supposés
    propres, $\opB\notin\varFil_i$. Par contradiction, on en déduit que
    $\varFilB$ est propre.
  \end{itemize}

  On applique donc le \cref{thm.Zorn}~: on trouve un élément maximal
  $\varUlF$ de $\varXX$ qui est supérieur à $\varFil$, c'est-à-dire
  qui contient $\varFil$. Il ne nous reste qu'à vérifier que $\varUlF$ est
  bien un ultrafiltre~: si $\varFilBB$ est un filtre propre contenant
  $\varUlF$, alors il contient $\varFil$, donc il appartient à
  $\varXX$. Par maximalité de $\varUlF$ dans $\varXX$, on en déduit
  que $\varUlF = \varFilBB$~: $\varUlF$ est un ultrafiltre sur $X$
  contenant $\varFil$.
\end{proof}

On peut maintenant montrer le \cref{thm.completion}~:

\begin{proof}
  Soit une signature $\Sign$ et une théorie $\varT$ sur $\Sign$ supposée
  cohérente, donc telle que pour toute partie finie $\varCG\subfin \varT$,
  $\bigwedge \varCG \nvdP \bot$. On redonne l'algèbre de Lindenbaum-Tarski~:
  \[\LinTar(\Sign) \defeq \quot{\Prop(\Sign)}{\dvvd}\]
  (comme on prouve ce théorème avant le théorème de complétude, on montre bien
  la proposition pour la relation $\vdP$ et non la relation $\vD$).

  La théorie $\varT$ induit sur $\LinTar(\Sign)$ une partie vérifiant
  la propriété d'intersection finie, et la théorie
  \[\vdclose{\varT} \defeq \compr{\varFF\in\Prop(\Sign)}{\varT \vdP \varFF}\]
  est ainsi un filtre propre~: on peut donc grâce au
  \cref{thm.ultrafilter.lemma} trouver un ultrafiltre $\varUlF$ contenant
  $\vdTclos$.

  Soit alors la théorie $\varT' = \bigcup\varUlF$, qui contient donc
  toutes les formules appartenant à une classe d'équivalence de $\varUlF$.
  On remarque que $\varT'$ est close par $\vdP$, donc puisque
  $\clas{\bot} \notin\varUlF$, cela signifie que $\bot\notin\varT'$.
  Il nous reste à vérifier que $\varT'$ est complète~: pour toute
  formule $\varFF$, on sait que $\varFF\lor\lnot\varphi \in \varT'$, donc
  $\clas{\varFF}\lor\clas{\lnot\varFF}\in\varUlF$, mais
  comme $\varUlF$ est un ultrafiltre, il est en particulier premier. On
  en déduit que soit $\clas\varFF \in \varUlF$, soit
  $\clas{\lnot\varFF}\in \varUlF$~: dans le premier cas,
  $\varFF\in \varT'$ et dans le deuxième cas,
  $\lnot\varFF\in \varT'$. Ainsi pour toute formule $\varFF$, soit
  $\varFF\in \varT'$ soit $\lnot\varFF\in \varT'$.

  Donc il existe une théorie complète $\varT'$ qui contient la théorie
  $\varT$.
\end{proof}

\noindent
On peut de plus étendre le filtre de Fréchet sur un ensemble en un
ultrafiltre.

\begin{proposition}
  Soit $X$ un ensemble infini. Alors il existe un ultrafiltre $\varUlF$ sur
  $\powerset(X)$ qui est différent de tous les $\FilPrinc{x}$. On appelle un
  ultrafiltre qui n'est pas principal un ultrafiltre libre.
\end{proposition}

\begin{proof}
  On sait que le filtre de Fréchet sur $X$ n'est inclus dans aucun des
  ultrafiltres $\FilPrinc{x}$, et d'après le lemme de l'ultrafiltre on peut
  trouver un ultrafiltre $\varUlF$ qui contient le filtre de Fréchet. Si
  $\varUlF = \FilPrinc{x}$ pour un certain $x\in X$, alors le filtre de
  Fréchet serait inclus dans ce $\FilPrinc{x}$, donc $\varUlF$ n'est pas
  principal.
\end{proof}

\subsection{Filtre et morphisme}

On définit en algèbre la notion de noyau d'une fonction comme l'ensemble des
éléments envoyés sur $0$. Les filtres, eux, peuvent être vus comme le
complémentaire d'un tel noyau~: dans le cas d'une algèbre de Boole, c'est même
littéralement le cas puisqu'un idéal sur une algèbre de Boole est le noyau d'un
morphisme d'anneau.

Nous allons donc voir comment des morphismes sur les différents treillis
permettent de définir de nouveaux filtres.

\begin{definition}[Rétrotirette d'un filtre]
  Soient $(X,\leq)$, $(Y,\leq)$ deux inf demi-treillis, $\varFil$ un filtre
  sur $Y$ et $\makeFunName{f}{X}{Y}$ une fonction croissante. On définit la
  rétrotirette de $\varFil$ par $f$ comme~:
  \[\pulba{f} \varFil \defeq \compr{x\in X}{f(x) \in \varFil}\]
\end{definition}

\begin{property}
  Soit $(X,\leq)$, $(Y,\leq)$ deux inf demi-treillis, $\varFil$ un filtre
  sur $Y$ et $f$ un morphisme d'inf demi-treillis, alors $\pulba{f}\varFil$
  est un filtre sur $X$.
\end{property}

\begin{proof}
  On vérifie les propriétés d'un filtre~:
  \begin{itemize}
  \item comme $f(\opT_X) = \opT_Y$ et $\opT_Y \in \varFil$, on en déduit que
    $\opT_X\in \pulba{f}\varFil$.
  \item si $x\in \pulba{f}\varFil$ et $y\in X$ tel que $x\leq y$, alors par
    on sait que $f(x)\in\varFil$, donc par croissance $f(x)\leq f(y)$ et,
    par clôture par le haut de $\varFil$, $f(y)\in \varFil$, donc
    $y\in \pulba{f}\varFil$.
  \item si $x\in \pulba{f}\varFil$ et $y\in \pulba{f}\varFil$, alors
    $f(x)\in\varFil$ et $f(y)\in\varFil$, donc
    $f(x)\opmt f(y)\in\varFil$ et comme $f$ est un morphisme,
    $f(x\opmt y)\in\varFil$. Donc $x\opmt y\in \pulba{f}\varFil$.\qedhere
  \end{itemize}
\end{proof}

Dans le cas d'un treillis, on peut aussi caractériser un filtre propre grâce à
l'existence de $\opB$.

\begin{property}
  Soient $(X,\leq)$ et $(Y,\leq)$ deux treillis, $\varFil$ un filtre propre
  sur $Y$ et $\makeFunName{f}{X}{Y}$ un morphisme de treillis. Alors
  $\pulba{f}\varFil$ est lui-même un filtre propre.
\end{property}

\begin{proof}
  Nous avons déjà vérifié que $\pulba{f}\varFil$ est un filtre. Il nous reste à
  montrer que celui-ci est propre. Mais comme on sait que $f(\opB_X)=\opB_Y$
  et que $\opB_Y\notin\varFil$, on en déduit que
  $\opB_X\notin \pulba{f}\varFil$, donc $\pulba{f}\varFil$ est propre.
\end{proof}

\noindent
Un morphisme de treillis préserve aussi le fait d'être premier.

\begin{property}
  Soient $(X,\leq)$ et $(Y,\leq)$ deux treillis, $\varFil$ un filtre premier
  sur $Y$ et $\makeFunName{f}{X}{Y}$ un morphisme de treillis. Alors
  $\pulba{f}\varFil$ est un filtre premier.
\end{property}

\begin{proof}
  Là encore, il nous suffit de montrer le caractère premier de
  $\pulba{f}\varFil$.
  Soient $x,y\in X$ tels que $x\opjn y \in \pulba{f}\varFil$, donc par
  définition
  $f(x\opjn y)\in \varFil$, mais comme $\varFil$ est premier et que
  $f(x\opjn y) = f(x)\opjn f(y)$, on en déduit que $f(x)\in\varFil$ ou
  $f(y)\in\varFil$, c'est-à-dire que $x\in \pulba{f}\varFil$ ou que
  $y\in \pulba{f}\varFil$. Donc $\pulba{f}\varFil$ est premier.
\end{proof}

Une conséquence directe est que si $\varFil$ est un filtre premier et que $X$
est une algèbre de Boole, $\pulba{f}\varFil$ est un ultrafiltre sur $X$.

D'un point de vue plus algébrique, si $\makeFunName{f}{\varBB}{\varBB'}$ est un
morphisme entre algèbres de Boole, alors le quotient de $\varBB$ par
$\ker(f)$ nous donne encore une algèbre de Boole (en tant que quotient
d'anneau), mais une propriété importante sur les anneaux est que le quotient par
un idéal maximal donne un corps. Dans le cas de nos algèbres de Boole, le
seul corps possible est $\quot{\bZ}{2\bZ}$, ce qui signifie qu'un ultrafiltre
correspond exactement au complémentaire du noyau d'un morphisme d'algèbre de
Boole vers $\quot{\bZ}{2\bZ}$.

\begin{proposition}
  Soit $(\varBB,\leq)$ une algèbre de Boole. Alors on a une bijection
  entre $\Hom(\varBB,\varBB_2)$, l'ensemble des morphismes de $\varBB$ vers
  $\quot{\bZ}{2\bZ}$, et l'ensemble $\Spec(\varBB)$ des ultrafiltres sur
  $\varBB$.
\end{proposition}

\begin{proof}
  Soit $\varUlF$ un ultrafiltre, on lui associe la fonction
  \[
  \makeFun{\indicP{\varUlF}}{\varBB}{\varBB_2}{x}
          {\begin{cases}
              \opT \text{ si } x \in \varUlF \\
              \opB \text{ si } x \notin \varUlF
          \end{cases}}\]

  Vérifions qu'on a bien un morphisme de treillis~:
  \begin{itemize}
  \item comme $\opT\in\varUlF$, $\indicP{\varUlF}(\opT) = \opT$, de même
    $\indicP{\varUlF}(\opB) = \opB$.
  \item pour montrer que
    $\indicP{\varUlF}(x\opmt y) = \indicP{\varUlF}(x)\opmt\indicP{\varUlF}(y)$,
    il suffit de montrer que
    \[x\opmt y \in \varUlF \iff x\in\varUlF \text{ et }y\in\varUlF\]
    Si $x\in \varUlF$ et $y\in\varUlF$, alors par compatibilité
    $x\opmt y \in \varUlF$. Si $x\opmt y \in \varUlF$, alors par clôture
    par le haut $x\in \varUlF$ et $y\in \varUlF$.
  \item pour montrer que
    $\indicP{\varUlF}(x\opjn y) = \indicP{\varUlF}(x)\opjn\indicP{\varUlF}(y)$,
    il suffit de montrer que
    \[x\opjn y \in \varUlF \iff x\in\varUlF \text{ ou }y\in\varUlF\]
    Si $x\in \varUlF$ (respectivement $y\in \varUlF$) alors comme
    $x\leq x\opjn y$ (respectivement $y\leq x\opjn y$) et par clôture par le
    haut,
    $x\opjn y \in \varUlF$. Si $x\opjn y \in \varUlF$, alors comme un
    ultrafiltre est premier, on en déduit que soit $x\in\varUlF$, soit
    $y\in \varUlF$.
  \end{itemize}

  Dans l'autre sens, étant donné un morphisme $\makeFunName{f}{\varBB}{\intB}$,
  on lui associe
  \[\varUlF_f \defeq \compr{x\in \varBB}{f(x) = \opT}\]
  On vérifie que $\varUlF_f$ est un ultrafiltre~:
  \begin{itemize}
  \item comme $f(\opT_B) = \opT$, $\opT\in \varUlF_f$.
  \item si $x\leq y$ et $f(x)=\opT$ alors $f(y)\geq \opT$ donc $f(y)=\opT$.
  \item si $f(x) = \opT$ et $f(y)=\opT$ alors
    $f(x)\opmt f(y) = \opT\opmt\opT =\opT$, et comme
    $f(x)\opmt f(y)=f(x\opmt y)$ on en déduit que $x\opmt y \in \varUlF_f$.
  \item $f(\opB_B)=\opB\neq\opT$ donc $f(\opB_B)\notin\varUlF_f$.
  \item si $x\opjn y \in \varUlF_f$ alors $f(x)\opjn f(y) = \opT$. Si
    $f(x) = \opB$ et $f(y) = \opB$ alors $f(x)\opjn f(y) \neq \opT$, donc par
    l'absurde on en déduit que $x\in \varUlF_f$ ou $y\in \varUlF_f$.
  \end{itemize}

  On vérifie aisément que $\varUlF \mapsto \indicP{\varUlF}$ et
  $f\mapsto \varUlF_f$ sont des fonctions réciproques l'une de l'autre.
  On a donc notre bijection.
\end{proof}

\begin{exercise}
  Soient $X,Y$ deux ensembles et $\makeFunName{f}{X}{Y}$ une fonction. On
  définit la fonction
  \[
  \makeFun{f^{-1}}{\powerset(Y)}{\powerset(X)}
          {\varPt}{\compr{x \in X}{f(x)\in \varPt}}
  \]
  Montrer que cette fonction est un morphisme d'algèbre de Boole.
\end{exercise}

\begin{notation}
  Pour une fonction $\makeFunName{f}{X}{Y}$ entre deux ensembles et un filtre
  $\varFil$ sur $\powerset(X)$, on notera $f\varFil$ la rétrotirette de
  $\varFil$ par $f^{-1}$. On n'a pas d'ambiguïté car, par défaut, $f$ n'est pas
  un morphisme d'algèbre de Boole entre les ensembles de parties.

  Donnons une description rapide de $f\varFil$~:
  \[f\varFil = \compr{\varPt \subseteq Y}{f^{-1}(\varPt) \in \varFil}\]
\end{notation}

\newpage

\section{Exercices et problèmes}

\subsection{Exercices}

\begin{exercise}[Sur l'implication]
  On se place dans une algèbre de Heyting $\varBH$.
  \begin{enumerate}[label=(\roman*)]
  \item Montrer que pour tous $a,x,y \in \varBH$, si $x \leq y$ alors
    $a \opimp x \leq a \opimp y$.
  \item Montrer que pour tous $x,y,a \in \varBH$, si $x \leq y$ alors
    $y \opimp a \leq x \opimp a$.
  \item En déduire que pour tous $x,y \in \varBH$, si $x \leq y$ alors
    $\opng y \leq \opng x$.
  \item Si $\varBH$ est aussi une algèbre de Boole, montrer que la
    réciproque de la question précédente est vraie aussi~: si
    $\opng y \leq \opng x$ alors $x \leq y$.
  \item Montrer que pour tous $x,y,z \in \varBH$,
    $x \opimp (y \opmt z) = (x \opimp y) \opmt (x \opimp z)$.
  \item Montrer que pour tous $x,y,z \in \varBH$,
    $(x\opjn y) \opimp z = (x \opimp z) \opjn (y \opimp z)$.
  \end{enumerate}
\end{exercise}

\begin{exercise}[Sur la négation]
  On se place dans une algèbre de Heyting $\varBH$.
  \begin{enumerate}[label=(\roman*)]
  \item Montrer que pour tout $x \in \varBH$,
    $\opng(\opng(\opng x)) = \opng x$.
  \item Montrer que pour tous $x,y \in \varBH$,
    $(x \to y) \leq (\opng y \opimp \opng x)$.
  \item Montrer que dans une algèbre de Boole, l'inégalité précédente est
    une égalité.
  \end{enumerate}
\end{exercise}

\begin{exercise}[Algèbre de Boole engendrée]
  Soit $(\varBB,\leq)$ une algèbre de Boole et $X\subseteq \varBB$.
  \begin{enumerate}[label=(\roman*)]
  \item Montrer qu'il existe une plus petite algèbre de Boole
    $\varBB(X)\subseteq \varBB$ contenant $X$.
  \end{enumerate}

  \noindent
  On définit une expression booléenne en $X$ par la grammaire
  \[e,e'\inddefeq x \mid \opng e \mid e\opjn e' \mid e \opmt e'\]
  où $x \in X$. On confond ici l'expression définie par induction et son
  interprétation évidente comme un élément de $\varBB$.
  \begin{enumerate}[label=(\roman*),resume]
  \item Montrer que $\varBB(X)$ est l'ensemble des expressions booléennes en
    $\varBB$.
  \end{enumerate}

  \noindent
  Pour tout ensemble $X$, on dit qu'une partie $\varPt \subseteq X$ est
  cofinie s'il existe un ensemble fini $F \subfin X$ tel que
  $\varPt=X\setminus F$. Montrer que l'ensemble des parties $\varPt\subseteq X$
  telles que $\varPt$ est finie ou cofinie forme une sous-algèbre de
  Boole de $\powerset(X)$.
  \begin{enumerate}[label=(\roman*),resume]
  \item Montrer que cette sous-algèbre de Boole est engendrée par
    l'ensemble $\compr{\{x\}}{x \in X}$.
  \end{enumerate}
\end{exercise}

\begin{exercise}[Atomes d'une algèbre de Boole]
  On appelle atome d'une algèbre de Boole $\varBB$ un élément
  $a \in \varBB$ tel que $a > \opB$ et tel qu'il n'existe pas d'élément
  $b \in \varBB$ tel que $\opB < b < a$.
  \begin{enumerate}[label=(\roman*)]
  \item Montrer que pour tout ensemble $X$ et élément $x \in X$,
    $\{x\}$ est un atome dans l'algèbre de Boole $\powerset(X)$.
  \end{enumerate}

  \noindent
  L'objectif de cet exercice est de montrer que toute algèbre de Boole
  finie $\varBB$ peut s'écrire comme un ensemble de parties. Pour ce faire, on
  va considérer l'ensemble $\At(\varBB)$ des atomes d'une algèbre de
  Boole.
  \begin{enumerate}[label=(\roman*),resume]
  \item Pour tout $x \in \varBB$, en notant
    \[\At(\varBB,x) \defeq \compr{y \in \At(\varBB)}{y < x}\]
    montrer que pour tout $x \in \varBB$, si $x \neq \opB$ alors
    $\At(\varBB,x) \neq \vide$.
  \item Montrer que $x = \opJn \At(\varBB,x)$.
    \textit{Indication~: en partant du fait que $x$ est un majorant de
      $\At(\varBB,x)$, montrer que l'égalité n'est pas stricte.}
  \item Montrer que
    \[
    \makeFun{f}{\varBB}{\powerset(\At(\varBB))}{x}{\At(\varBB,x)}
    \]
    est un isomorphisme d'algèbres de Boole.
  \item Une algèbre de Boole infinie est-elle encore isomorphe à
    l'ensemble des parties d'un ensemble~?
    \textit{Indication~: trouver une algèbre de Boole infinie sans
      atome.}
  \end{enumerate}
\end{exercise}

\begin{exercise}[Quotient d'algèbres de Heyting]
  Dans le cas des algèbre de Boole, le complémentaire d'un filtre étant
  un idéal, on possède une façon naturelle de construire une algèbre de
  Boole quotient. \'Etudions le cas d'une algèbre de Heyting.

  Soit $\varBH$ une algèbre de Heyting et $\varFil$ un filtre sur
  $\varBH$. On définit la relation $\sim_{\varFil}$ par
  \[x\sim_{\varFil} y \defeq (x \opimp y) \opmt (y\opimp x) \in \varFil\]
  \begin{enumerate}[label=(\roman*)]
  \item Montrer que $\sim_{\varFil}$ est une relation d'équivalence.
  \item Montrer que si $x \sim_{\varFil} x' $ et $y \sim_{\varFil} y'$, alors
    $x\opjn y \sim_{\varFil} x'\opjn y'$, $x\opmt y \sim_{\varFil} x'\opmt y$
    et $x\opimp y \sim_{\varFil} x' \opimp y'$.
  \item Vérifier que $\quot{\varBH}{\sim_{\varFil}}$ est bien une algèbre de
    Heyting avec les opérations induites sur celles de $\varBH$.
  \item Montrer que si $\varBB$ est une algèbre de Boole et $\varFil$ est
    un filtre sur $\varBB$, alors en posant
    \[\varId \defeq \compr{\varBB \setminus \varPt}{\varPt \in \varFil}\]
    on a pour tous $x,y \in \varBB$ l'équivalence
    \[x \sim_{\varFil} y \iff x \opXOR y \in \varId\]
  \end{enumerate}

  \noindent
  On peut aussi construire un analogue du \ordinalnumeralmale{1} théorème
  d'isomorphisme. Pour cela, la première étape est de définir la notion de
  morphisme adaptée aux algèbres de Heyting. Pour deux algèbres de
  Heyting $\varBH,\varBH'$, on dit que $\makeFunName{f}{\varBH}{\varBH'}$
  est un morphisme lorsque $f$ est un morphisme de treillis qui, de plus,
  commute avec l'implication~:
  \[\forall x,y \in \varBH, f(x \opimp y) = f(x) \opimp f(y)\]
  Dans le reste de l'exercice, on fixe $\varBH,\varBH'$ deux algèbres de
  Heyting et $\makeFunName{f}{\varBH}{\varBH'}$ un morphisme entre les
  deux.
  \begin{enumerate}[label=(\roman*),resume]
    \item Montrer que l'ensemble
    \[\ker(f) \defeq \compr{x \in \varBH}{f(x) = \opT}\]
    est un filtre sur $\varBH$.
  \item Montrer que si $\varBH_0\subseteq \varBH$ est une sous-algèbre de
    Heyting de $\varBH$, alors $f(\varBH_0)$ (l'image directe de
    $\varBH_0$ par $f$) est une sous-algèbre de Heyting de $\varBH'$. En
    déduire que $\im(f)$ est une algèbre de Heyting. Montrer de même que
    si $\varBH'_0\subseteq\varBH'$ est une sous-algèbre de Heyting de
    $\varBH'$, alors $f^{-1}(\varBH'_0)$ est une sous-algèbre de Heyting
    de $\varBH$.
  \item Montrer que la fonction induite par $f$
    \[\makeFun{\relev{f}}{\quot{\varBH}{\ker(f)}}{\im(f)}{\clas{x}}{f(x)}\]
    est bien définie et est un isomorphisme d'algèbres de Heyting.
    \textit{Note~: le quotient par l'ensemble $\ker(f)$ est à lire comme
      le quotient par la relation associée au filtre $\ker(f)$.}
  \end{enumerate}
\end{exercise}

\newpage

\subsection{Problèmes}

%\begin{problem}[Une structure algébrique pour le calcul des prédicats]
%  Nous avons vu qu'il est possible d'axiomatiser grâce à la théorie des ordres
%  la structure de la logique propositionnelle (à équivalence logique près).
%  L'objectif de cet exercice est d'étendre cette idée en donnant des structures
%  qui se comportent comme l'ensemble des formules du calcul des prédicats, en
%  ajoutant en particulier des propriétés que vérifient les quantificateurs et
%  l'égalité.
%  
%  On se fixe dans cet exercise une signature $\Sign$. Pour tout
%  $n \in \bN$, on dénote par $\Formula_n(\Sign)$ l'ensemble des formules
%  de $\Sign$ dont les variables libres appartiennent à
%  $\{x_0,\ldots,x_{n-1}\}$ (ainsi $\Prop(\Sign) = \Formula_0(\Sign)$).
%
%  On définit aussi $\LinTar_n(\Sign)$, l'algèbre de Lindenbaum-Tarski des
%  formules à $n$ variables libres, comme le quotient de $\Formula_n(\Sign)$
%  par la relation $\dvvd$.
%  \begin{enumerate}[label=(\roman*)]
%  \item Montrer que chaque $\LinTar_n(\Sign)$ est une algèbre de Boole.
%  \item Soit $\iota_n : \Formula_n(\Sign) \to \Formula_{n+1}(\Sign)$
%    l'injection canonique des formules à $n$ variables libres dans les formules
%    à $n+1$ variables libres. Justifier que $\iota_n$ est compatible avec
%    $\dvvd$. On note alors aussi $\iota_n$ la fonction canoniquement
%    associée $\LinTar_n(\Sign) \to \LinTar_{n+1}(\Sign)$~: montrer que
%    cette fonction est un morphisme d'algèbre de Boole.
%  \item De même, montrer que
%    \[
%    \makeFun{\mathord{\forall}}{\Formula_{n+1}(\Sign)}{\Formula_n(\Sign)}
%            {\varFF}{\forall x_n,\varFF}
%    \]
%    induit une fonction croissante entre $\LinTar_{n+1}(\Sign)$ et
%    $\LinTar_n(\Sign)$. Faire de même avec le quantificateur
%    $\mathord{\exists}$.
%  \item Montrer que pour tous $i,j < n$, la fonction
%    \[
%    \makeFunName{\kron{i}{j}}{\Formula_{n+1}(\Sign)}{\Formula_n(\Sign)}
%    \]
%    définie en substituant $x_i$ par $x_j$ et chaque $x_{k+1}$ par $x_k$ pour
%    $k > i$ (pour assurer que les variables sont dans $\{x_0,\ldots,x_n\}$),
%    induit un morphisme d'algèbre de Boole entre $\LinTar_{n+1}(\Sign)$ et
%    $\LinTar_n(\Sign)$.
%  \item Montrer l'équivalence suivante pour
%    $\varFF \in \LinTar_n(\Sign)$ et
%    $\varFP \in \LinTar_{n+1}(\Sign)$~:
%    \[\varFF \vdP \forall \varFP \iff \iota_n(\varFF)\vdP \varFP\]
%    ainsi que, pour $\varFF \in \LinTar_{n+1}(\Sign)$ et
%    $\varFP \in \LinTar_n(\Sign)$~:
%    \[\exists \varFF \vdP \varFP \iff \varFF\vdP \iota_n(\varFP)\]
%    Ces équivalences caractérisent en un sens les quantificateurs $\forall$ et
%    $\exists$.
%  \item Montrer l'équivalence suivante, pour $\varFF \in \LinTar_n(\Sign)$
%    et $\varFP \in \LinTar_{n+1}(\Sign)$~:
%    \[\iota_n(\varFF) \land x_i = x_j \vdP \varFP \iff
%    \varFF\vdP \kron{i}{j}(\varFP)\]
%    Cet équivalence caractérise en un sens l'égalité.
%  \end{enumerate}
%\end{problem}

\begin{problem}[Logique temporelle linéaire \cite{LTL77}]
  Ce problème étudie une variante de logique appelée logique temporelle linéaire
  (\foreignexpr{Linear Time Logic} en anglais, abrégé en $\LTL$).
  On la verra dans le cas d'une logique propositionnelle~: son objectif est de
  traduire des propositions qui peuvent évoluer dans le temps.

  \pbPart{Logique modale linéaire}

  Pour commencer cette étude, on travaille dans une logique simplifiée, qu'on
  appelle Logique Modale Linéaire ($\LML$). Dans cette logique, les propositions
  ne sont pas à interpréter comme étant une valeur de vérité
  \og Vrai\fg{} / \og Faux\fg mais comme une valeur de vérité évoluant dans le
  temps~: la sémantique d'une proposition sera un prédicat sur l'ensemble des
  fonctions $\bN \to \intB$.
  On notera $\infWd$ l'ensemble des fonctions $\bN \to \intB$ (cette notation
  aura plus de sens à partir du \cref{chp.ordinaux}).

  Donnons la syntaxe des propositions de $\LML$. L'ensemble $\Prop_{\LML}$ est
  construit inductivement par la grammaire suivante~:
  \[\varFF,\varFP\inddefeq x \mid \vara \mid \top\mid \bot \mid \varFF\land
  \varFP \mid \varFF\lor \varFP \mid \lnot \varFF\mid \modNext\varFF\]
  où $x \in \Var$. La constante $\vara$ est utilisée comme une proposition
  atomique. On définit récursivement, pour un environnement d'évaluation
  $\env : \Var \partialto \powerset(\infWd)$ et une proposition
  $\varFF\in\Prop_{\LML}$ telle que $\VL(\varFF) \subseteq \dom(\env)$, la
  valuation $\ValE(\varFF)$~:
  \begin{itemize}
  \item $\ValE(x) = \env(x)$
  \item $\ValE(\vara) = \compr{\varstr \in \infWd}{\varstr(0) = 1}$
  \item $\ValE(\top) = \powerset(\infWd)$
  \item $\ValE(\bot) = \vide$
  \item $\ValE(\varFF \land \varFP) = \ValE(\varFF) \cap \ValE(\varFP)$
  \item $\ValE(\varFF \lor  \varFP) = \ValE(\varFF) \cup \ValE(\varFP)$
  \item $\ValE(\lnot\varFF) = \powerset(\infWd)\setminus \ValE(\varFF)$
  \item $\ValE(\modNext\varFF)=
    \compr{\varstr\in\infWd}{\upstr\in\ValE(\varFF)}$
  \end{itemize}
  Où $\upstrPP{\varstr}{k}$ est la fonction $n \mapsto \varstr(n+k)$, qui est
  donc la suite $\varstr$ décalée vers la gauche.
  Si $\varFF$ est close, on note directement
  $\Val(\varFF)$ pour la valeur commune à tous les $\ValE(\varFF)$.

  \begin{pbQuest}
  \item Soit une suite binaire finie $b \defeq b_0\ldots b_{n-1} \in \intB^n$.
    Construire une formule $\varFF_b$ telle que
    $\Val(\varFF_b) = \compr{\varstr\in\infWd}{\varstr\restr n = b}$.
  \end{pbQuest}

  \noindent
  On définit l'équivalence de formules par
  \[\varFF\equiv\varFP \defeq \forall \env\supset \VL(\varFF)\cup\VL(\varFP),
  \ValE(\varFF) = \ValE(\varFP)\]
  \begin{pbQuest}[resume]
  \item Montrer que l'algèbre de Lindenbaum-Tarski
    \[\LinTar(\LML) = \quot{\Prop_{\LML}}{\equiv}\]
    est une algèbre de Boole, avec les opérations induites par les
    constructeurs des formules.
  \item Montrer que $\modNext$ induit la fonction
    \[
    \makeFun{\modNext}{\LinTar(\LML)}{\LinTar(\LML)}
            {\clas{\varFF}}{\clas{\modNext\varFF}}
    \]
    qui est un morphisme d'algèbres de Boole.
  \end{pbQuest}

  \noindent
  On désigne par $\LinTar_0(\LML)$ le sous-ensemble de $\LinTar(\LML)$
  constitué des classes de propositions sans variables.
  \begin{pbQuest}[resume]
  \item Montrer que $\LinTar_0(\LML)$ est une sous-algèbre de Boole de
    $\LinTar(\LML)$.
  \item Montrer que $\Val$ est un morphisme d'algèbre de Boole de
    $\LinTar_0(\LML)$ dans $\powerset(\infWd)$.
  \end{pbQuest}

  \noindent
  On souhaite maintenant montrer que $\Val$ ne constitue pas un isomorphisme.
  \begin{pbQuest}[resume]
  \item Montrer que $\Val$ est, en tant que fonction de domaine
    $\LinTar_0(\LML)$, une injection.
  \end{pbQuest}

  \noindent
  On cherche donc à prouver que $\Val$ n'est pas surjective. Pour cela, le
  moyen le plus immédiat est de trouver des éléments de $\powerset(\infWd)$
  sans antécédent.
  \begin{pbQuest}[resume]
  \item Montrer que pour toute formule $\varFF$ ne contenant pas de variable,
    $\Val(\varFF)$ est un ensemble infini ou $\vide$.
  \item En déduire que $\{0^\infty\}$ contenant uniquement la suite constante
    nulle, n'est pas définissable dans $\LML$.
  \end{pbQuest}

  \pbPart{La logique $\LTL$}

  On considère maintenant la logique $\LTL$, dans laquelle on ajoute
  l'opérateur binaire $\modU$, dont on définit la sémantique par
  \[\ValE(\varFF\modU\varFP) \defeq
  \compr{\varstr\in\infWd}{\exists i \in \bN,
    \begin{array}{l}
      \upstrPP{\varstr}{i} \in\ValE(\varFF)\\
      \forall j < i, \upstrPP{\varstr}{j}\in\ValE(\varFP)
    \end{array}
  }\]
  On appelle \og $\varFF$ jusqu'à $\varFP$\fg la proposition
  $\varFF\modU\varFP$.
  L'ensemble des propositions $\Prop_{\LTL}$ est alors
  \[\varFF,\varFP\inddefeq\cdots\mid \varFF\modU\varFP\]

  \begin{pbQuest}
  \item Montrer les équivalences suivantes pour toutes propositions
    $\varFF,\varFP,\varFC$~:
    \begin{center}
      \begin{tabular}{cc}
        $\modNext(\varFF\modU\varFP) \equiv
        (\modNext\varFF)\modU(\modNext\varFP)$
        &
        $\varFF\modU(\varFP\lor\varFC)\equiv
        (\varFF\modU\varFP)\lor(\varFF\modU\varFP)$
        \\
        \\
        \multicolumn{2}{c}{
          $(\varFF\land\varFP)\land\varFC\equiv
          (\varFF\modU\varFC)\lor(\varFP\modU\varFC)$}
      \end{tabular}
    \end{center}
  \item On définit $\modR$, le dual de De Morgan de $\modU$, par
    \[\varFF\modR\varFP\defeq \lnot(\lnot\varFF\modU\lnot\varFP)\]
    Décrire $\ValE(\varFF\modR\varFP)$.
  \end{pbQuest}

  \noindent
  On souhaite d'abord proposer deux autres modalités, dont on va montrer
  qu'elles peuvent être définies dans $\LTL$. On pose la modalité
  \og un jour \fg et la modalité \og toujours \fg, notées respectivement
  $\modDia\varFF$ et $\modSq\varFF$ avec la sémantique
  \begin{align*}
    \ValE(\modDia\varFF) &\defeq
    \compr{\varstr\in\infWd}
           {\exists i\in\bN,\upstrPP{\varstr}{i}\in\ValE(\varFF)}
    \\
    \ValE(\modSq\varFF) &\defeq
    \compr{\varstr\in\infWd}
           {\forall i\in\bN,\upstrPP{\varstr}{i}\in\ValE(\varFF)}
  \end{align*}
  \begin{pbQuest}[resume]
  \item Pour une proposition $\varFF\in\Prop_{\LTL}$, construire une proposition
    $\varFF'$ équivalente à $\modDia\varFF$, c'est-à-dire telle que
    $\ValE(\varFF') = \ValE(\modDia\varFF)$. Faire de même avec la modalité
    $\modSq$.
  \item Montrer que pour toute proposition $\varFF$, on a les équivalences
    \[\modDia\varFF \equiv \varFF\lor\modNext(\modDia\varFF)\qquad
    \modSq\varFF \equiv \varFF\land\modNext(\modSq\varFF)\]
  \item Montrer pour toute proposition les deux équivalences
    \[\modDia\varFF \equiv \lnot(\modSq(\lnot\varFF))
    \modSq\varFF \equiv \lnot(\modDia(\lnot\varFF))\]
  \item Quel est le sens de la proposition $\modSq\modDia\varFF$~? Et de la
    proposition $\modDia\modSq\varFF$~?
  \item Montrer qu'il existe une proposition $\varFF\in\Prop_{\LTL}$ telle que
    pour toute proposition $\varFP\in\Prop_{\LML}$, $\varFF\not\equiv\varFP$.
  \end{pbQuest}

  \pbPart{$\LTL$ et la logique du premier ordre}

  Dans cette partie, on donne une autre présentation d'un langage pour parler de
  propriétés au cours du temps. On considère ici la logique du premier ordre,
  telle que décrite dans le \cref{chp.logpred}. On fixe la signature
  \[\Sign \defeq \{<^2, \vara^1\}\]
  où $<$ et $\vara$ sont des symboles de relation. Pour une suite
  $\varstr\in\infWd$, on définit le modèle $\varMM_{\varstr}$ par
  \begin{itemize}
  \item L'ensemble $\carM{\varMM_{\varstr}}$ est l'ensemble des couples
    $(n,\varstr(n))$, c'est-à-dire le graphe de $\varstr$ vu comme une
    fonction.
  \item La relation $(n,\varstr(n)) < (m,\varstr(m))$ est donnée par
    $n < m$.
  \item La relation $\vara((n,\varstr(n)))$ est donnée par
    $\varstr(n) = 1$.
  \end{itemize}

  \noindent
  Pour une suite $\varstr\in\infWd$, on définit $\varstr\satisf\varFF$ par
  $\varMM_{\varstr}\satisf\varFF$.

  \begin{pbQuest}
  \item Construire une formule $\Ss(x,y)$ telle que
    \[\varstr,(x\mapsto (n,\varstr(n)),y\mapsto(m,\varstr(m)))\satisf\Ss(x,y)
    \iff m = n + 1\]
  \end{pbQuest}

  \noindent
  On définit la sémantique d'une proposition $\varFF\in\Prop(\Sign)$ par
  \[\semFOL{\varFF}\defeq \compr{\varstr\in \infWd}{\varstr\satisf\varFF}\]
  \begin{pbQuest}
  \item Construire une proposition $\varFF$ telle que $\semFOL{\varFF}$ est
    l'ensemble des suite commençant par $0$.
  \item \'Etant donnée une proposition $\varFF$, construire une proposition
    $\modNext\varFF$ telle que
    \[\semFOL{\modNext\varFF} \defeq
    \compr{\varstr\in\infWd}{\upstr\satisf\varFF}\]
  \item Montrer que
    $\semFOL{\varFF\lor\varFP} = \semFOL{\varFF}\lor\semFOL{\varFP}$,
    $\semFOL{\varFF\land\varFP} = \semFOL{\varFF}\land\semFOL{\varFP}$ et
    $\semFOL{\lnot\varFF} = \powerset(\infWd)\setminus \semFOL{\varFF}$
  \item En déduire que pour toute proposition $\varFF\in\Prop_{\LML}$ ne
    contenant pas de variable, il existe $\varFP \in \Prop(\Sign)$ telle que
    $\semFOL{\varFP} = \Val(\varFF)$.
  \item \'Etendre le résultat précédent à $\LTL$.
  \end{pbQuest}
\end{problem}
